\chapter{管路、孔口、管嘴的水力计算}

\begin{yijubox}
E1 \texttt{8-1考情分析.pptx} slide33：第五章 管路、孔口、管嘴的水力计算＝\textbf{考试计算题重点}。\\
E2 \texttt{9-1考点分析.pptx}：计算题占 50\% 分值；第四章是主要考点、需掌握计算套路；\textbf{第五章同为主要计算考点}。\\
E1 slide16 各题型考法：简答考\emph{水头损失分类及定义}；选择判断题考\emph{概念辨析}；填空考\emph{名词定义与小计算}；计算考\emph{伯努利方程、动量方程}。\\
E4 真题落点：2022 年 818 单选第 3、7、8、9、10 题均在本章；2015 年试题填空考"层流断面平均速度＝0.5 倍最大速度""$\Rey$ 判流态""U 形水银压差计"，计算第 3 题考"突然缩小水头损失"；2026 回忆版简答考"水击现象""层流与雷诺数关系"，计算第 5 题考"管嘴体积流量（书上例题）"；题库 \texttt{10.pdf} 选择第 10 题、计算第 2 题也在本章。\\
\textbf{本章结论}：这是全书\textbf{篇幅最大、计算题命中率最高}的一章。一句话概括考法——
选择填空考\emph{判据与系数}（雷诺数、长管/短管、光滑/粗糙、串联/并联），简答考\emph{水头损失分类与水击}，
计算题考\emph{长管/短管 + 孔口/管嘴}。复习策略＝把判据背成条件反射，把五六个系数（$\lambda$、$\zeta$、$\mu$、$\varphi$、$\varepsilon$、$c$）的取值与含义记牢。
\end{yijubox}

\section{流动阻力和水头损失}

\begin{kaodian}{5}{水头损失的两种形式与叠加原理}
\textbf{是什么：}实际流体都有黏性，流动过程中流体内部、流体与管壁之间都存在摩擦与撞击，
机械能不可逆地转化为热能，沿流动方向总水头一定下降，这部分损失叫\textbf{水头损失}，用 $h_w$ 表示。

\textbf{两种形式（必背，考简答/填空）：}
\begin{center}
\begin{tabularx}{\textwidth}{@{}lXX@{}}
\toprule
类别 & 成因 & 发生位置 \\
\midrule
沿程水头损失 $h_f$ & 黏性摩擦\emph{沿整个流程连续作用}（质点之间、流体与管壁之间） & 等直径直管段 \\
局部水头损失 $h_j$ & 边界急剧变化使流速的大小与方向被迫重组，产生\textbf{漩涡}，漩涡耗能 & 弯头、阀门、变径、进出口等局部管件 \\
\bottomrule
\end{tabularx}
\end{center}

阻力还可按内因、外因分：流体之间的摩擦与掺混＝\textbf{内部阻力}（受管径、流量、黏度影响）；
流体与管壁的摩擦撞击＝\textbf{外部阻力}（受接触面积、管壁粗糙度、流量影响）。

\noindent\textbf{教材原话（p69）：}「在边界无变化（流动的方向、壁面的性质、过流断面的形状和尺寸等均不变）的均匀流段内，
由于黏性形成的阻碍流体运动的力称为沿程阻力，或称摩擦阻力。」「在边界发生急变的流段内，由于过流断面变化、
流动方向改变，速度重新分布，质点间进行动量交换而产生的阻力称为局部阻力。」
教材并注明「通常渐变流段产生的阻力也视为沿程阻力」，$\lambda$ 的取值「与流体的流动状态和壁面特性有关，由实验确定」。

\begin{gongshi}{水头损失的叠加原理（教材式 5.4）}
\[ h_w=\sum h_f+\sum h_j \]
\textbf{含义：}一段管路的总水头损失＝各直管段沿程水头损失之和＋各局部管件的局部水头损失之和。
之所以能直接相加，是因为它们消耗的都是同一股流体的机械能，且各自独立计算。
教材式（5.1）$p_w=\rho g h_w$ 即这条叠加原理的压强表达（$p_w$ 为单位体积流体的压强损失，教材 p69）。\\
\textbf{适用条件：}恒定流、液体充满全管；两个局部损失之间隔有足够长的直管段，彼此不干扰。
\end{gongshi}

\textbf{常考题型：}简答（"水头损失分几类？各自定义与成因？"）；填空（"总水头损失＝\rule{2cm}{0.4pt}"）；
选择（判断某一损失属沿程还是局部）。
\textbf{重要度 \xstarfull{5}}\quad\yiju{E1 slide16「简答考水头损失分类及定义」；E1 slide33「第五章＝考试计算题重点」；E2「第五章同为主要计算考点」}\quad\nandu{易}

\begin{banfa}
记忆锚点：\textbf{"沿程＝黏性沿程摩擦，局部＝漩涡耗能"}。
考场上一看到"弯头、阀门、突然扩大、出口"就写 $h_j$；一看到"管长 $l$、管径 $d$"就写 $h_f$。
简答答题模板：先分两类，再各给一句成因，最后写 $h_w=\sum h_f+\sum h_j$ 收尾。
\end{banfa}
\end{kaodian}

\begin{kaodian}{4}{有效断面的水力要素：湿周、水力半径与当量直径}
\textbf{是什么：}决定沿程阻力大小的断面几何量。流道面积 $A$ 越大，内部阻力越小；
湿周 $\chi$（有效断面上流体与固体边壁接触的周界长度）越大、管壁越粗糙，外部阻力越大。

\noindent\textbf{出处说明：}湿周、水力半径与当量直径的定义见教材 3.2.7 节「湿周、水力半径和当量直径」（书 p36），
第 5 章用到该结论时直接引用，本节内容与该定义一致。

\begin{gongshi}{湿周、水力半径、当量直径与相对粗糙度}
\[ \chi\ (\mathrm{m}),\qquad R_h=\frac{A}{\chi}\ (\mathrm{m}),\qquad d_e=4R_h,\qquad \frac{\Delta}{d} \]
\textbf{含义：}$R_h=A/\chi$ 称\textbf{水力半径}，综合反映断面的"通畅"程度——$R_h$ 越大，阻力越小；
$d_e=4R_h$ 称\textbf{当量直径}，作用是把圆管的一切公式直接搬到非圆断面上；
$\Delta$ 为管壁粗糙突起的平均高度（\textbf{绝对粗糙度}），$\Delta/d$ 为\textbf{相对粗糙度}。\\
\textbf{适用条件：}当量直径只用于\textbf{雷诺数与沿程阻力}的折算，\emph{不能}用于过流面积与流量计算，
即不等于 $\pi d_e^2/4$（圆管除外）。
\end{gongshi}

\textbf{常用断面的水力半径：}圆管 $R_h=\dfrac{\pi d^2/4}{\pi d}=\dfrac{d}{4}$；
矩形 $a\times b$：$R_h=\dfrac{ab}{2(a+b)}$；同心环缝：$R_h=\dfrac{D-d}{4}$。

\textbf{常用管壁当量粗糙度 $\Delta$（教材表 5.1，p78）：}铜管、新氯乙烯管及玻璃管 0.001$\sim$0.002 mm；
钢管 0.03$\sim$0.07 mm；涂锌铁管 0.1$\sim$0.2 mm；新铸铁管 0.2$\sim$0.4 mm；旧铸铁管 0.5$\sim$1.5 mm；
焊接钢管（中度生锈）0.5 mm；混凝土管 0.3$\sim$3.0 mm。

\begin{tishi}
\textbf{订正（以教材为准）：}原稿此处列的是外校辅导书口径的绝对粗糙度（普通钢管 0.19 mm、干净铸铁管 0.45 mm、
污秽金属管 0.75$\sim$0.9 mm、水管道 0.25$\sim$1.25 mm、混凝土槽 0.8$\sim$9.0 mm）。
教材表 5.1 给的是\textbf{当量粗糙度}：钢管 0.03$\sim$0.07 mm、新铸铁管 0.2$\sim$0.4 mm、
旧铸铁管 0.5$\sim$1.5 mm、混凝土管 0.3$\sim$3.0 mm——与原稿数值可差数倍乃至一个量级。
考研以教材为准，故正文改用教材表 5.1 的数值；两说的差别来源是：原稿用的是实测绝对粗糙度的推荐值，
教材用的是「阻力效果与人工粗糙相当」的\textbf{当量粗糙度}（教材 p78 定义）。
\end{tishi}

\textbf{常考题型：}填空（给矩形或环形断面求 $R_h$、$d_e$）；选择（"当量直径能否用于求流量"）；
计算（非圆管算雷诺数）。
\textbf{重要度 \xstarfull{4}}\quad\yiju{E5 题库选择第 10 题（圆管与方管沿程损失比 0.886，必须先算当量直径）；E1 slide33「第五章＝考试计算题重点」}\quad\nandu{中}

\begin{yicuo}
① \textbf{当量直径不是"等效面积"}：它是 $d_e=4A/\chi$。若误按 $d_e=\sqrt{4A/\pi}$ 再代入 $Q=Av$，答案必错。
② 只有圆管才有 $d_e=d$；矩形、环缝都不等，别把特征尺寸当 $d_e$。
③ 水力半径 $R_h$ 与几何半径 $r_0$ 差 2 倍：圆管 $R_h=r_0/2=d/4$。
④ "湿周"必须是与\textbf{固体边壁}接触的那部分周界，明渠的\emph{自由水面不算湿周}。
\end{yicuo}
\end{kaodian}

\section{黏性流体的两种流态}

\begin{kaodian}{5}{雷诺实验与两种流态}
\textbf{实验（雷诺实验）：}英国物理学家雷诺用一只稳水位箱、一支水平玻璃管和一根颜色水针做实验。
步骤与现象：保持水箱水位恒定，微开出口阀门，颜色水在玻璃管内呈\emph{一条清晰的直线}；
逐渐开大阀门，颜色水线开始\emph{抖动、弯曲}；开度足够大后颜色水\emph{扩散、布满全管}。
反向关小阀门可倒序复现。

\textbf{两种流态（定义必背）：}
\begin{center}
\begin{tabularx}{\textwidth}{@{}lXX@{}}
\toprule
流态 & 特征 & 阻力来源 \\
\midrule
层流 & 流体质点\emph{互不干扰、各自成层}，沿管轴做平滑直线运动，层间不混合不碰撞 & 层间黏性摩擦力 \\
湍流（紊流） & 流体质点\emph{互相掺混、杂乱无章}；主体做轴向运动的同时有\textbf{径向脉动} & 黏性摩擦＋掺混与涡旋 \\
\bottomrule
\end{tabularx}
\end{center}

\textbf{过渡流}只是层流到湍流的中间不稳定状态，\textbf{不是独立流型}（流型只有层流与湍流两种）；
工程计算中常把它归入湍流。

\noindent\textbf{教材原话（p70）：}「流体质点做规律的线状运动，彼此互不混掺的运动称层流」；
「如果流体质点在运动中出现不规则的互相混掺，质点运动方向随机变化，这种流动称为湍流」。
教材并指出是「英国科学家雷诺在 1883 年给出了判定两种流态的准则」。

\noindent\textbf{教材原话（p71）：}「在实际工程中，扰动普遍存在，上临界流速没有实际意义，一般的临界流速指的是下临界流速」。

\textbf{临界流速：}流态转变时的流速。层流转为湍流时称\textbf{上临界流速} $v_{cu}$；
湍流转为层流时称\textbf{下临界流速} $v_{cd}$，且 $v_{cu}>v_{cd}$（说明流态转变有"迟滞"）。

\textbf{水头损失的实验测定：}在等直径水平试验段两端接出两根测压管，稳定流时由伯努利方程得
$h_f=\dfrac{p_1-p_2}{\rho g}$，据此画出 $h_f$ 与 $v$ 的关系曲线，即得两种流态的不同规律。

\begin{jielun}
\textbf{层流：}$h_f\propto v^{1.0}$（直线）；\textbf{湍流：}$h_f\propto v^{1.75\sim2.0}$（曲线）。
这两段"直线／曲线"就是雷诺实验最重要的定量结论，简答与选择题都爱考。
\end{jielun}

\textbf{常考题型：}简答（描述雷诺实验现象与结论）；选择（"过渡流是否为独立流型"）；
空（"层流时 $h_f$ 与 $v$ 的几次方成正比"）。
\textbf{重要度 \xstarfull{5}}\quad\yiju{E4 2026 回忆版简答考"层流与雷诺数关系"；E1 slide16「简答考水头损失分类及定义」；E1 slide33「第五章＝考试计算题重点」}\quad\nandu{易}

\begin{banfa}
实验现象用\textbf{三句话}背：直线 $\to$ 抖动 $\to$ 满管扩散。
被问"为什么两种流态的水头损失规律不同"，落脚点是：\emph{层流只有黏性摩擦，湍流多了质点掺混与涡旋，耗能随速度增长更快}。
上、下临界流速的关系记成"上去难、下来容易"，即 $v_{cu}>v_{cd}$。
\end{banfa}
\end{kaodian}

\begin{kaodian}{5}{层流与湍流的判别标准——雷诺数}
\textbf{是什么：}实验证明，影响流态的因素是流速 $v$、管径 $d$、密度 $\rho$、黏度 $\mu$ 四个量；
它们组成的复合数群 $\dfrac{\rho v d}{\mu}$ 就是判别流态的准则，称\textbf{雷诺数}。
$v$、$d$、$\rho$ 越大、$\mu$ 越小，越容易从层流转变为湍流。

\begin{gongshi}{雷诺数与流态判据（教材式 5.5、5.6、5.7）}
\[ \Rey=\frac{\rho v d}{\mu}=\frac{vd}{\nu}\quad\text{（圆管）},\qquad
   \Rey=\frac{v d_e}{\nu}=\frac{4vR_h}{\nu}\quad\text{（非圆管）} \]
\[ \Rey_c\approx 2300\ \text{（下临界）；工程判据：}\ \Rey<2000\ \text{层流},\quad \Rey>2300\ \text{湍流} \]
\textbf{含义：}$\Rey$ 的物理意义是\textbf{惯性力与黏性力之比}。
$\Rey$ 大 $\to$ 惯性力占主导 $\to$ 湍流；$\Rey$ 小 $\to$ 黏性力占主导 $\to$ 层流。
所以 $\Rey$ 越大，湍流程度越高。\\
\textbf{适用条件：}特征长度必须与特征速度配对——圆管取 $d$，非圆管取当量直径 $d_e$；
$\nu$、$\mu$ 取该温度下的值。\textbf{下临界雷诺数才是判别标准}（上临界雷诺数不稳定、随扰动变化大，不能作判据）。
\end{gongshi}

\textbf{临界雷诺数的取值：}实验测得下临界 $\Rey_c\approx 2300$，上临界可高到 13800；
工程计算为安全起见常把临界值取为 \textbf{2000}——$\Rey<2000$ 按层流算，$\Rey>2000$ 按湍流算。

\begin{tishi}
\textbf{订正（判据口径，以教材为准）：}原稿把 2000 写成唯一判据。教材第 5 章（p72）给出的判据是
$\Rey<2320$ 为层流、$\Rey>2320$ 为湍流；$\Rey<2000$ 层流、$\Rey>2000$ 湍流是工程上更保守的常用口径
（历年真题答案即按 2000 给出）。两说并列：答案落在 $2000\sim2320$ 之间时，按题目给定的临界值作答；
题目未给时写明所用口径，\textbf{同一题内不要两套混用}。
\end{tishi}

\begin{zhenti}{2022 年 818 单选第 9 题}{单项选择题}
直径分别为 $d=10\ \mathrm{mm}$、$d=30\ \mathrm{mm}$ 的两根圆管，管内水流运动黏度相同、平均流速相同时，判别两管中水流的流态。
\end{zhenti}
\noindent\textbf{要点：}两管流速相同，$\Rey=\dfrac{vd}{\nu}\propto d$，管径大的雷诺数大。分别算出 $\Rey_1$、$\Rey_2$ 后与 2000（或 2300）比较，得到"小管层流、大管湍流"这类结论。\textbf{关键是把 $d$ 换算成米。}

\begin{zhenti}{2015 年试题 填空}{填空题（原题）}
水管直径为 $10\ \mathrm{mm}$，管中水的平均流速 $v=0.2\ \mathrm{m/s}$，其运动黏度 $\nu=1.308\times10^{-6}\ \mathrm{m^2/s}$，
则管中水流的流态为\rule{2cm}{0.4pt}；如将管径改为 $30\ \mathrm{mm}$、流速不变时，流态为\rule{2cm}{0.4pt}。
\end{zhenti}
\noindent\textbf{解：}$\Rey_1=\dfrac{0.2\times0.01}{1.308\times10^{-6}}=1529<2000$，\textbf{层流}；
$\Rey_2=\dfrac{0.2\times0.03}{1.308\times10^{-6}}=4587>2000$，\textbf{湍流}。

\textbf{常考题型：}选择/填空（给 $v$、$d$、$\nu$ 判流态，或问 $\Rey$ 的物理意义）；
简答（"层流与雷诺数的关系"）。
\textbf{重要度 \xstarfull{5}}\quad\yiju{E4 2022 年 818 单选第 9 题；E4 2015 年试题填空"$\Rey$ 判流态"；E4 2026 回忆版简答"层流与雷诺数关系"；E1 slide33「第五章＝考试计算题重点」}\quad\nandu{中}

\begin{banfa}
\textbf{判流态三步走：}① 统一单位（mm$\to$m）；② $\Rey=\dfrac{vd}{\nu}$（给的是动力黏度 $\mu$ 就先算 $\nu=\mu/\rho$）；
③ 与 2000/2300 比。\\
\textbf{反推式：}题目问"欲保持层流，最大流速是多少"$\to$ 令 $\Rey=2000$，得 $v_{\max}=\dfrac{2000\nu}{d}$。
题目问"最小管径能达到多大流量仍为层流"同理令 $\Rey=2000$。\\
\textbf{速记：}$\nu$ 的量纲是 $\mathrm{L^2/T}$，$\Rey$ 无量纲；$\Rey$ 中的 $d$ 一定是\emph{内径}。
\end{banfa}

\begin{yicuo}
① \textbf{判据用下临界雷诺数}，不用上临界雷诺数。
② $\Rey=\dfrac{vd}{\nu}$ 用的是\textbf{运动黏度} $\nu$；给动力黏度 $\mu$ 时必须先除密度，这一步历年都在丢分。
③ 非圆管必须用\textbf{当量直径} $d_e=4R_h$，不能随便取某个边长。
④ 层流与湍流的 $h_f\sim v$ 关系是 $v^1$ 与 $v^{1.75\sim2}$，不要写成 $v$ 与 $v^2$ 以外的东西。
\end{yicuo}
\end{kaodian}

\section{圆管中的层流流动}

\begin{kaodian}{4}{圆管层流的切应力分布与速度分布}
\textbf{推导思路（理解即可，不要求证）：}在半径为 $r_0$ 的圆管中取半径 $r$、长 $l$ 的圆柱体液流，
两端压差 $\Delta p=p_1-p_2$ 推着它走，周围表面的黏性摩擦力拖着它，等速流动时二者平衡，
再用牛顿内摩擦定律积分即得速度分布。

\begin{gongshi}{圆管层流的切应力分布、速度分布（教材式 5.8、5.9、5.10、5.15）}
\[ \tau=\frac{\rho g J r}{2}=\frac{\Delta p}{2l}r,\qquad
   u=\frac{\rho g J}{4\mu}\left(r_0^2-r^2\right)=\frac{\Delta p}{4\mu l}\left(r_0^2-r^2\right) \]
\textbf{含义：}切应力 $\tau$ 在断面上呈\textbf{线性分布}——管轴处为零，管壁处最大
$\tau_0=\dfrac{\rho g J r_0}{2}$；速度 $u$ 沿半径呈\textbf{抛物线（旋转抛物面）分布}——
管轴处最大、管壁处为零。\\
\textbf{适用条件：}只适用于\textbf{圆管、层流、恒定流、不可压缩牛顿流体}；
$J=h_f/l$ 为水力坡度（单位长度的水头损失）。
\end{gongshi}

\textbf{速度分布的五个特点（简答常考）：}① 抛物线形；② 管中心流速最大；③ 沿壁方向渐减；
④ 管壁处流速为零（黏附条件）；⑤ 断面平均流速为最大流速的一半。

\textbf{常考题型：}选择（"圆管层流过流断面上切应力的分布"$\to$管轴处为零、与半径成正比）；
简答（层流速度分布特点）。
\textbf{重要度 \xstarfull{4}}\quad\yiju{E4 2015 年试题选择（圆管层流过流断面切应力分布）；E4 2015 年试题填空"层流断面平均速度＝0.5 倍最大速度"；E1 slide33「第五章＝考试计算题重点」}\quad\nandu{中}

\begin{yicuo}
① 切应力分布是\textbf{线性}、速度分布是\textbf{抛物线}，两者千万不要互换。
② 切应力在\emph{管轴为零、管壁最大}；速度在\emph{管轴最大、管壁为零}——方向正好相反。
③ 水力坡度 $J=h_f/l$ 只对沿程损失成立，它与总水头线的坡度不是一回事，别混用。
\end{yicuo}
\end{kaodian}

\begin{kaodian}{5}{圆管层流的流量公式与沿程阻力系数}
\begin{gongshi}{最大流速、平均流速、哈根--泊肃叶流量公式（教材式 5.11、5.12、5.13、5.14）}
\[ u_{\max}=\frac{\rho g J r_0^2}{4\mu}=\frac{\Delta p\,r_0^2}{4\mu l},\qquad
   v=\frac{u_{\max}}{2}=\frac{\Delta p\,r_0^2}{8\mu l}=\frac{\Delta p\,d^2}{32\mu l} \]
\[ Q=\frac{\pi r_0^4\Delta p}{8\mu l}=\frac{\pi d^4\Delta p}{128\mu l} \]
\textbf{含义：}层流时管轴处流速最大；\textbf{断面平均流速恰为最大流速的一半}
（这一条是历年填空的常客）；流量与\textbf{管径的四次方}成正比，与压差成正比、与管长和黏度成反比。
哈根--泊肃叶公式还可写成 $Q=\dfrac{\pi d^4\Delta p}{128\mu l}$，医学、化工里常用它测黏度。\\
\textbf{适用条件：}圆管、\textbf{层流}、恒定流、充分发展的不可压缩牛顿流体；
管径越小、黏度越大，越容易满足层流条件。
\end{gongshi}

\begin{gongshi}{层流的沿程水头损失与沿程阻力系数 $\lambda$（教材式 5.2、5.16、5.17）}
\[ h_f=\frac{\Delta p}{\rho g}=\frac{32\mu l v}{\rho g d^2},\qquad
   \lambda=\frac{64}{\Rey}=\frac{64\mu}{\rho v d},\qquad
   h_f=\lambda\frac{l}{d}\frac{v^2}{2g} \]
\textbf{含义：}层流时沿程水头损失与平均流速的\textbf{一次方}成正比（与雷诺实验一致）；
把上式改写成达西公式的形式，就得到层流的沿程阻力系数 $\lambda=\dfrac{64}{\Rey}$——
它\emph{只与雷诺数有关，与管壁粗糙度无关}。\\
\textbf{适用条件：}\textbf{仅适用于圆管层流}（$\Rey<2300$）；非圆管层流要用 $d_e$ 并把系数改为 $C/\Rey$
（如矩形断面 $C\approx 57$）。湍流绝对不能套 $64/\Rey$。
\end{gongshi}

\begin{zhenti}{2022 年 818 单选第 10 题}{单项选择题}
层流状态下，若管径减为原来的一半（流量不变、管长与流体性质不变），沿程压强损失增为原来的多少倍？
\end{zhenti}
\noindent\textbf{要点（答 16 倍）：}由 $h_f=\dfrac{32\mu lv}{\rho g d^2}$ 且 $v=\dfrac{4Q}{\pi d^2}$ 得
$h_f\propto\dfrac{v}{d^2}\propto\dfrac{1}{d^4}$。$d\to d/2$ 时 $h_f$ 变为 $2^4=16$ 倍。
\textbf{记住这条结论：层流、$Q$ 不变时 $h_f\propto d^{-4}$。}

\textbf{常考题型：}计算（求 $Q$、$v$、$\Delta p$、$\mu$ 之一）；选择（"管径减半损失变为几倍"）；
填空（"层流平均流速＝最大流速的\rule{1.2cm}{0.4pt}倍"）。
\textbf{重要度 \xstarfull{5}}\quad\yiju{E4 2022 年 818 单选第 10 题（层流管径减半、压强损失倍数，答 16 倍）；E4 2015 年试题填空"层流断面平均速度＝0.5 倍最大速度"；E1 slide33「第五章＝考试计算题重点」}\quad\nandu{中}

\begin{banfa}
\textbf{层流算题固定套路：}① 判流态（不判就做风险极大）；② 层流则 $\lambda=64/\Rey$；
③ 需要压差就用 $h_f=\dfrac{\Delta p}{\rho g}$ 或用 $Q=\dfrac{\pi d^4\Delta p}{128\mu l}$；④ 检查单位。\\
\textbf{比例速算表（层流、$Q$ 不变）：}$h_f\propto d^{-4}$；$h_f\propto l^{1}$；$h_f\propto\mu^{1}$。
\textbf{层流、$v$ 不变：}$h_f\propto d^{-2}$。做题先问自己"题里固定的是 $Q$ 还是 $v$"。
\end{banfa}
\end{kaodian}

\section{湍流的沿程水头损失}

\begin{kaodian}{3}{湍流的脉动与时均化、混合长度}
\textbf{是什么：}湍流中任一点的速度大小、方向都随时间做\emph{高频无规则}的随机变化，这叫\textbf{脉动}。
脉动使湍流不能用确定性的瞬时速度描述，工程上采用\textbf{时均化}：把瞬时速度写成时均值加脉动值。

\begin{gongshi}{湍流的时均速度与脉动值（教材式 5.18、5.19、5.20、5.21）}
\[ \overline{u}=\frac{1}{T}\int_0^T u\,\dd t,\qquad u=\overline{u}+u',\qquad \overline{u'}=0 \]
\textbf{含义：}取足够长的时段 $T$，把瞬时速度对时间取平均即得\textbf{时均速度} $\overline{u}$；
瞬时值与其之差为\textbf{脉动值} $u'$，脉动量的时均值恒为零。
工程上所说的"湍流速度"通常指时均速度。\\
\textbf{适用条件：}时间平均的时段 $T$ 相对于脉动周期足够长、相对于流动变化的宏观特征时间又足够短（时均恒定）。
\end{gongshi}

\textbf{脉动带来的附加切应力（雷诺应力）：}脉动使流体微团在层间交换动量，产生附加切应力
$\tau_t=-\rho\overline{u'v'}$；加上黏性切应力 $\mu\dfrac{\dd\overline{u}}{\dd y}$ 就是湍流中的总切应力。
\textbf{混合长度}是普朗特为描述这种掺混而引入的假想距离，常写作 $l$，反映微团在失去原有动量前走过的距离。

\textbf{常考题型：}选择/判断（"脉动是湍流的基本特征"、"时均值为零的是脉动值"）。
\textbf{重要度 \xstarfull{3}}\quad\yiju{E1 slide33 只要求"概念与公式套用"；E2「第五章同为主要计算考点」（本点为概念铺垫，无独立大题）}\quad\nandu{易}

\begin{banfa}
一句话辨脉动：\textbf{"湍流＝时均流动＋脉动"}；\textbf{"脉动量的时均值恒为零"}。
选择题出现"湍流的速度分布是抛物线"$\to$错（湍流是较\emph{丰满}的指数型剖面，抛物线是层流）。
\end{banfa}
\end{kaodian}

\begin{kaodian}{4}{层流底层、水力光滑与水力粗糙}
\textbf{是什么：}湍流并非整个断面都乱——紧贴管壁仍存在一层极薄的\textbf{黏性底层（层流底层）}，
厚度记 $\delta$；$\delta$ 与管壁粗糙突起高度 $\Delta$ 的大小关系，决定了管壁粗糙度对流动有没有影响。

\noindent\textbf{教材原话（p76）：}层流底层「厚度很薄，通常不到 1mm」；「如果 $\delta>\Delta$，管壁的粗糙突起将全部淹没在
层流底层之中，核心湍流好像在一完全光滑管道中流动，这种管道称为水力光滑管」；「当 $\delta<\Delta$，管壁的粗糙突起
大部暴露在层流底层之外……这种管道称为水力粗糙管」。

\begin{gongshi}{层流底层厚度与光滑/粗糙判别式（教材式 5.22）}
\[ \delta\approx\frac{32.8\,d}{\Rey\sqrt{\lambda}} \]
\[ \delta>\Delta\ \Rightarrow\ \text{水力光滑管};\qquad \delta<\Delta\ \Rightarrow\ \text{水力粗糙管} \]
\textbf{含义：}当层流底层完全淹没管壁的粗糙突起（$\delta>\Delta$）时，湍流区感受不到粗糙度的影响，
如同在完全光滑的管中流动，称\textbf{水力光滑管}，此时 $\lambda$ 只与 $\Rey$ 有关；
反之粗糙突起暴露在湍流区（$\delta<\Delta$），流体撞上突起会生成漩涡、新增能量损失，称\textbf{水力粗糙管}，
此时 $\lambda$ 与相对粗糙度有关。\\
\textbf{适用条件：}判别用 $\delta$ 与 $\Delta$ 的比较，工程上常用经验判据 $\Rey\sqrt{\lambda}$ 与 $\dfrac{\Delta}{d}$ 的关系。
同一根管子流量变大时 $\delta$ 变薄，可能从"水力光滑"变成"水力粗糙"。
\end{gongshi}

\textbf{常考题型：}选择（"水力光滑管是指管壁绝对粗糙度为零"$\to$错）；
简答（"什么叫水力光滑与水力粗糙"）；填空（层流底层越薄越接近粗糙管）。
\textbf{重要度 \xstarfull{4}}\quad\yiju{E1 slide33「第五章＝考试计算题重点」；E2「第五章同为主要计算考点」}\quad\nandu{中}

\begin{yicuo}
① \textbf{"水力光滑"不等于"几何光滑"}：它是由层流底层厚度与 $\Delta$ 的相对大小决定的\emph{流动状态}，
同一根管子在小流量下是水力光滑、大流量下是水力粗糙。
② 层流底层\textbf{越厚}越光滑，$\delta$ 随 $\Rey$ 增大而\emph{减薄}。
③ 不要用"新管/旧管"这类几何说法直接判光滑粗糙，必须先算 $\Rey$ 与 $\Delta/d$。
\end{yicuo}
\end{kaodian}

\begin{kaodian}{5}{尼古拉兹实验分区、莫迪图与 $\lambda$ 的计算}
\textbf{是什么：}尼古拉兹用人工均匀砂粒粗糙管做系统实验，测出沿程阻力系数 $\lambda$ 与
$\Rey$、相对粗糙度 $\Delta/d$ 的关系，把流动分成五个区——这就是本章最"硬"的考点。

\noindent\textbf{教材原话（p78）：}「当量粗糙度是指阻力效果与人工粗糙相当的绝对粗糙度」——
工业管道按此把实测阻力折算成表 5.1 的当量粗糙度 $\Delta$ 值（教材表 5.1，p78）

\begin{gongshi}{达西公式（沿程水头损失通用式，教材式 5.2）}
\[ h_f=\lambda\frac{l}{d}\frac{v^2}{2g}=\lambda\frac{8Q^2l}{\pi^2gd^5}=0.0826\,\lambda\frac{Q^2l}{d^5} \]
\textbf{含义：}$\lambda$ 称沿程阻力系数（水力摩擦系数），它与 $\Rey$、$\Delta/d$ 有关。
用 $Q$ 表达的形式 $0.0826\lambda Q^2l/d^5$ 在管路计算中最好用，因为它把 $v$ 消掉了。\\
\textbf{适用条件：}\textbf{圆管、恒定流}；层流与湍流都适用，区别只在 $\lambda$ 取哪个公式；
非圆管要用 $d_e$ 替换 $d$。
\end{gongshi}

\textbf{尼古拉兹五个分区（必背）与 $\lambda$ 公式：}
\begin{center}
\begin{tabularx}{\textwidth}{@{}lXX@{}}
\toprule
分区 & 范围 & $\lambda$ 的计算 \\
\midrule
I 层流区 & $\Rey<2320$ & $\lambda=\dfrac{64}{\Rey}$，与 $\Delta/d$ 无关 \\
II 临界过渡区 & $2320<\Rey<4000$ & $\lambda$ 随 $\Rey$ 急剧增大，规律不稳定，工程上按湍流处理 \\
III 湍流光滑区 & $4000<\Rey<26.98\left(\dfrac{d}{\Delta}\right)^{8/7}$ & $\lambda$ 只与 $\Rey$ 有关：布拉修斯式 $\lambda=\dfrac{0.3164}{\Rey^{0.25}}$（$4000<\Rey<10^5$） \\
IV 湍流过渡区 & $26.98\left(\dfrac{d}{\Delta}\right)^{8/7}<\Rey<4160\left(\dfrac{d}{2\Delta}\right)^{0.85}$ & $\lambda=f\left(\Rey,\dfrac{\Delta}{d}\right)$，用柯罗布鲁克公式 \\
V 湍流粗糙区 & $\Rey>4160\left(\dfrac{d}{2\Delta}\right)^{0.85}$ & $\lambda$ 只与 $\dfrac{\Delta}{d}$ 有关，$h_f\propto v^2$，故又称\textbf{阻力平方区} \\
\bottomrule
\end{tabularx}
\end{center}

\begin{gongshi}{分区实用公式（教材式 5.23、5.24、5.26、5.27）}
\[ \text{光滑区：}\ \frac{1}{\sqrt{\lambda}}=2\lg\left(\Rey\sqrt{\lambda}\right)-0.8 \]
\[ \text{过渡区（柯罗布鲁克）：}\ \frac{1}{\sqrt{\lambda}}=-2\lg\left(\frac{\Delta}{3.7d}+\frac{2.51}{\Rey\sqrt{\lambda}}\right) \]
\[ \text{粗糙区（尼古拉兹）：}\ \frac{1}{\sqrt{\lambda}}=1.14-2\lg\frac{\Delta}{d},\qquad
   \text{或（希弗林松）：}\ \lambda=0.11\left(\frac{\Delta}{d}\right)^{0.25} \]
\textbf{湍流水力光滑区的两个实用式（教材式 5.23、5.24）：}
\[ \lambda=\frac{0.3164}{\Rey^{0.25}}\quad(4000<\Rey<10^5,\ \text{布拉修斯式}) \]
\[ \lambda=0.0032+\frac{0.221}{\Rey^{0.237}}\quad(10^5<\Rey<10^6,\ \text{尼古拉兹式}) \]
\textbf{含义：}光滑区 $\lambda$ 只随 $\Rey$ 减小而减小；过渡区 $\lambda$ 同时受 $\Rey$ 与 $\Delta/d$ 控制；
粗糙区 $\lambda$ 与 $\Rey$ 无关，\emph{只}取决于相对粗糙度，此时 $h_f\propto v^2$。\\
\textbf{适用条件：}教材过渡区除柯罗布鲁克式外还给出洛巴耶夫式（教材式 5.25，其式子在扫描件上残缺、
无法判读，本书不列）；柯罗布鲁克式是过渡区的通用式；粗糙区公式在 $4160(d/2\Delta)^{0.85}$ 以上使用；
人工粗糙管实验结论外推到工业管道时要偏保守（工业管按"等当量粗糙度"折算）。
\end{gongshi}

\begin{tishi}
\textbf{订正（术语，以教材为准）：}原稿把 Colebrook 公式写作「柯列勃洛克」，教材（p79）写作「柯罗布鲁克」，
并说明莫迪图「湍流水力过渡区是按柯罗布鲁克式绘制的」（教材 p79）。全章已统一为教材用语。
\end{tishi}

\begin{gongshi}{沿程水头损失的经验公式：谢才公式与曼宁公式（教材式 5.28~5.31）}
\[ v=C\sqrt{RJ},\qquad h_f=\frac{l\,v^2}{C^2R},\qquad C=\frac{1}{n}R^{1/6} \]
\textbf{含义：}$C$ 称\textbf{谢才系数}、$R$ 为水力半径、$J=h_f/l$ 为水力坡度。教材指出谢才公式是法国工程师
谢才 1775 年归纳的经验式（教材式 5.28）；将式（5.29）与达西公式结合，可得 $C=\sqrt{8g/\lambda}$（教材 p79）。
$C$ 常用两条经验式计算：曼宁公式 $C=R^{1/6}/n$（教材式 5.30），或巴甫洛夫斯基公式
$C=\dfrac{1}{n}R^y$，其中 $y=2.5\sqrt{n}-0.13-0.75\sqrt{R}\left(\sqrt{n}-0.10\right)$，
工程上近似取 $y=1.5\sqrt{n}$（$R<1.0\ \mathrm{m}$）或 $y=1.3\sqrt{n}$（$R>1.0\ \mathrm{m}$）（教材式 5.31）。\\
\textbf{适用条件：}教材（p79）特别指出\textbf{曼宁公式与巴甫洛夫斯基公式「只适用于湍流水力粗糙区」}；
巴甫洛夫斯基式的适用范围为 $0.1\le R\le3.0\ \mathrm{m}$、$0.011\le n\le0.04$；
粗糙系数 $n$ 按渠道材料与养护情况查教材表 5.2（p79--80）。
\end{gongshi}

\textbf{莫迪图：}工程上把 $\lambda$、$\Rey$、$\Delta/d$ 的关系绘成\textbf{莫迪图}，横坐标 $\Rey$、纵坐标 $\lambda$、
以 $\Delta/d$ 为参数。用法：算出 $\Rey$ 与 $\Delta/d$，查图得 $\lambda$。图中最上方一条水平线就是粗糙区。

\begin{zhenti}{题库 \texttt{10.pdf} 选择题第 10 题}{单项选择题}
圆管和正方形管道的断面面积、长度、相对粗糙度都相等，流动属湍流粗糙区，通过的流量相等，
则两种形状管道的沿程水头损失之比为（\quad）。\\
(A) 1.255\quad (B) 0.886\quad (C) 3.125\quad (D) 0.258
\end{zhenti}
\noindent\textbf{要点（答 B）：}粗糙区 $\lambda$ 只与 $\Delta/d$ 有关，而两者相对粗糙度相等且面积相等，
故 $\lambda$ 相同；由 $h_f\propto lQ^2/(\lambda^{-1}d^5)$ 归结为 $h_f\propto d^{-5}$。
圆管 $d$ 与方管当量直径 $d_e=4A/\chi$ 之比决定结果，得 $h_{\text{圆}}/h_{\text{方}}=0.886$。
\textbf{关键动作：方管必须换算当量直径 $d_e=4A/\chi$。}

\textbf{常考题型：}选择/填空（判断某流动属哪个区、给 $\Rey$ 与 $\Delta/d$ 求 $\lambda$）；
计算（先用 $\lambda$ 公式再算 $h_f$，常需\emph{迭代}）。
\textbf{重要度 \xstarfull{5}}\quad\yiju{E5 题库选择第 10 题（圆管与方管沿程损失比 0.886）；E1 slide33「第五章＝考试计算题重点」；E2「第五章同为主要计算考点」}\quad\nandu{难}

\begin{banfa}
\textbf{判区口诀：}先算 $\Rey$——小于 2320 层流（$64/\Rey$）；再算 $26.98(d/\Delta)^{8/7}$ 与 $4160(d/2\Delta)^{0.85}$
两个界值，就落进光滑区／过渡区／粗糙区。
\textbf{工程速判：}相对粗糙度小时常在光滑区用布拉修斯式；相对粗糙度大、$\Rey$ 大时按粗糙区取
$\lambda=0.11(\Delta/d)^{0.25}$；拿不准就一律用柯罗布鲁克式（通用）。
\textbf{迭代技巧：}$\lambda$ 出现在等式两边时，先取 $\lambda_0=0.03$ 代进去算两次就能收敛，
比"试差法"更快；若用迭代，写清初值与迭代次数，阅卷可给过程分。
\end{banfa}

\begin{yicuo}
① \textbf{$\lambda=64/\Rey$ 只适用于圆管层流}，湍流套用必错。
② 布拉修斯式 $\lambda=0.3164/\Rey^{0.25}$ 只适用于\textbf{湍流光滑区}（约 $4000<\Rey<10^5$）。
③ 粗糙区 $\lambda$ \textbf{与 $\Rey$ 无关}，只看 $\Delta/d$，不要再去查 $\Rey$。
④ 临界过渡区规律不稳定，考试若遇到直接按湍流处理，不要试图精确算 $\lambda$。
⑤ 莫迪图读数有误差，最后一步务必核对数量级（水管的 $\lambda$ 一般在 0.02$\sim$0.05 之间）。
\end{yicuo}
\end{kaodian}

\begin{liti}{5.1}{书 p80--81 · 沿程损失·判区与达西公式}
\textbf{题目}（教材原题，p80）内径 $d=0.2\ \mathrm{mm}$ 的钢管输送水，流量 $Q=0.04\ \mathrm{m^2/s}$，
水的运动黏度 $v=1.007\times10^{-6}\ \mathrm{m^2/s}$，钢管内壁的绝对粗糙度 $\Delta=0.04\ \mathrm{mm}$。
求 $1000\ \mathrm{m}$ 管道上的沿程水头损失 $h_f$。\\
\textbf{题目照录与两说并列：}$d=0.2\ \mathrm{mm}$、$Q=0.04\ \mathrm{m^2/s}$、$v=1.007\times10^{-6}\ \mathrm{m^2/s}$
是教材题干的印刷原文；而教材本题\emph{下文}是按 $d=0.2\ \mathrm{m}$、$Q=0.04\ \mathrm{m^3/s}$、
$\nu=1.007\times10^{-6}\ \mathrm{m^2/s}$ 计算的。两说并列，本讲义\textbf{不替教材改圆}，
计算以数字关系为准（判定见盒内末尾的说明）。

\begin{jie}
\textbf{① 求断面平均流速。}按与全部数字自洽的那一说，取 $d=0.2\ \mathrm{m}$，则
\[ A=\frac{\pi d^2}{4}=\frac{\pi\times0.2^2}{4}=0.0314\ \mathrm{m^2} \]
\[ v=\frac{Q}{A}=\frac{0.04}{0.1^2\pi}=1.27\ \mathrm{m/s} \]
\textbf{② 判别流态。}
\[ \Rey=\frac{vd}{\nu}=\frac{1.27\times0.2}{1.007\times10^{-6}}=2.52\times10^5>2320 \]
因 $\Rey>2320$，流动\textbf{不是层流}，须继续判别湍流分区。\\
\textbf{③ 判别湍流分区。}相对粗糙度为
\[ \frac{d}{\Delta}=\frac{0.2}{0.04\times10^{-3}}=5000 \]
水力光滑区的上界为
\[ 26.98\left(\frac{d}{\Delta}\right)^{8/7}=26.98\times5000^{8/7}\approx4.56\times10^5 \]
因 $4000<\Rey=2.52\times10^5<4.56\times10^5$，故流动处于\textbf{湍流水力光滑区}，
$\lambda$ 应按尼古拉兹式（教材式 5.24）计算。\\
\textbf{④ 计算 $\lambda$。}
\[ \lambda=0.0032+\frac{0.221}{\Rey^{0.237}} \]
\[ \lambda=0.0032+\frac{0.221}{\left(2.52\times10^5\right)^{0.237}}=0.0032+0.0116=0.0148 \]
\textbf{⑤ 由达西公式求 $h_f$。}
\[
\begin{aligned}
h_f&=\lambda\frac{l}{d}\frac{v^2}{2g}\\
   &=0.0148\times\frac{1000}{0.2}\times\frac{1.27^2}{2\times9.8}=6.09\ \mathrm{m}
\end{aligned}
\]
即单位重量水流经 $1000\ \mathrm{m}$ 管道的沿程损失为\textbf{6.09 m 水柱}。
\end{jie}

\begin{tishi}
\textbf{教材例 5.1 的两处印误（p80）：}① 题干印作「内径 $d=0.2\ \mathrm{mm}$」，但同一题下文按
$v=\dfrac{Q}{A}=\dfrac{0.04}{0.1^2\pi}=1.27\ \mathrm{m/s}$ 计算，由 $A=\pi d^2/4=0.0314\ \mathrm{m^2}$ 反推 $d=0.2\ \mathrm{m}$，
故题干中的「mm」应为「m」；② 题干把 $Q$ 的单位印成 $\mathrm{m^2/s}$（应为 $\mathrm{m^3/s}$），
「运动黏度」又写成 $v$（应为 $\nu$）。做题时以数字关系为准，按 $d=0.2\ \mathrm{m}$ 计。
\end{tishi}
\end{liti}

\begin{liti}{5.2}{书 p81 · 明渠·谢才公式与曼宁公式}
\textbf{题目}（教材原题，p81）设有一梯形断面的土渠，如图 5.13 所示。已知养护情况中等，
渠道底宽 $b=5\ \mathrm{m}$、水深 $h=2.5\ \mathrm{m}$、边坡系数 $m=\cot\alpha=1$。试求谢才系数 $C$ 值。

\tuyi{梯形断面土渠：底宽 $b=5\ \mathrm{m}$、水深 $h=2.5\ \mathrm{m}$，两侧边坡与水平面成 $\alpha$，边坡系数 $m=\cot\alpha=1$。}

\begin{jie}
\textbf{① 求过流断面面积与湿周。}梯形断面的几何关系为
\[ A=bh+mh^2,\qquad \chi=b+2h\sqrt{1+m^2} \]
代入数据：
\[ A=bh+mh^2=5\times2.5+1\times2.5^2=18.75\ \mathrm{m^2} \]
\[ \chi=b+2h\sqrt{1+m^2}=5+2\times2.5\times\sqrt{1+1^2}=12.07\ \mathrm{m} \]
（湿周 $\chi$ 只算与土渠\emph{固体边壁}接触的那部分周界，\textbf{自由水面不计入}。）\\
\textbf{② 求水力半径。}
\[ R=\frac{A}{\chi}=\frac{18.75}{12.07}=1.55\ \mathrm{m} \]
\textbf{③ 查粗糙系数。}由教材表 5.2（土渠、养护情况中等）查得 $n=0.0225$。\\
\textbf{④ 按曼宁公式（教材式 5.30）计算谢才系数。}
\[ C=\frac{1}{n}R^{1/6}=\frac{1}{0.0225}\times1.55^{1/6}=47.81\ \mathrm{m^{1/2}/s} \]
\textbf{⑤ 按巴甫洛夫斯基公式（教材式 5.31）计算。}指数为
\[
\begin{aligned}
y&=2.5\sqrt{n}-0.13-0.75\sqrt{R}\left(\sqrt{n}-0.10\right)\\
 &=2.5\times0.15-0.13-0.75\times1.2450\times(0.15-0.10)\\
 &=0.375-0.13-0.047=0.20
\end{aligned}
\]
代入巴甫洛夫斯基式：
\[ C=\frac{1}{n}R^{y}=\frac{1}{0.0225}\times1.55^{0.20}=48.51\ \mathrm{m^{1/2}/s} \]
两种经验公式所得 $C$ 相差约 $1.5\%$，工程上均可采用。教材把两个结果都列出，正说明
$C$ 是\emph{经验值}：同一断面用不同经验式会有小幅差异，$C$ 的量纲为 $\mathrm{m^{1/2}/s}$。
\end{jie}

\begin{tishi}
\textbf{考纲属性（重要）：明渠，818 考纲外，仅作了解。}
本题是\textbf{梯形断面土渠}（明渠）的谢才系数计算，属明渠均匀流内容，对应旧大纲／复试范围的
\textbf{考点码 X1（明渠流与堰流）}，现行 818 初试\textbf{不考}——旧大纲真题（2015 年第 10 题、2016 年、
2018 年计算题）确实反复考过「梯形断面土渠水力最佳断面」，但那些都按 X1 处理。
不过下面两条仍然要记：① 谢才公式 $v=C\sqrt{RJ}$ 与达西公式的联系 $C=\sqrt{8g/\lambda}$（教材 p79），
只是换了符号，物理上是同一件事；② 梯形断面的几何量 $A=bh+mh^2$、$\chi=b+2h\sqrt{1+m^2}$、$R=A/\chi$
在第 3 章「湿周与水力半径」考点里就要用到。明渠的具体算例本讲义不再扩展。
\end{tishi}
\end{liti}

\section{局部水头损失}

\begin{kaodian}{4}{局部水头损失公式与突然扩大（波达公式）}
\textbf{是什么：}流体经过局部管件时，边界形状突变使主流与壁面脱离、形成漩涡，
流速重新分布、质点互相撞击，由此产生的机械能损失即\textbf{局部水头损失} $h_j$。

\begin{gongshi}{局部水头损失通用式（教材式 5.3）}
\[ h_j=\zeta\frac{v^2}{2g} \]
\textbf{含义：}$\zeta$ 称\textbf{局部阻力系数}，无量纲，由管件几何形状决定；
$v$ 一般取\emph{管件下游}（较小断面）的断面平均流速——$\zeta$ 的数值必须与所用 $v$ 配套，
这是查表时最容易出错的一点。\\
\textbf{适用条件：}局部损失与沿程损失互不干扰（管件之间隔有足够长的直管）；
一般用实验测定的 $\zeta$，几何相似时 $\zeta$ 在湍流粗糙区近似为常数。
\end{gongshi}

\begin{gongshi}{突然扩大的局部水头损失（波达公式，教材式 5.35~5.37）}
\[ h_j=\frac{\left(v_1-v_2\right)^2}{2g},\qquad
   \zeta_1=\left(1-\frac{A_1}{A_2}\right)^2,\qquad
   \zeta_2=\left(\frac{A_2}{A_1}-1\right)^2 \]
\textbf{含义：}突然扩大管中，局部损失等于\textbf{两断面流速差的动能}——
原因是大断面一侧出现环形漩涡区，损失与速度差直接挂钩。
$\zeta_1$ 对应小断面流速 $v_1$，$\zeta_2$ 对应大断面流速 $v_2$。\\
\textbf{教材的推导路线（教材式 5.32~5.34，p82）：}先对突扩前后两断面写伯努利方程（式 5.32），
再写动量方程（式 5.34，因突扩段壁面压强分布未知，用动量方程绕过），联立即得式（5.35）波达公式。
教材指出「对一些特殊情况，如圆管突然扩大的流动则可以用理论方法导出」，其余局部损失只能实验测定。\\
\textbf{适用条件：}湍流（该式在层流中不成立）；若 $A_2\to\infty$（管出口流入大容器），
则 $h_j=\dfrac{v_1^2}{2g}$，即 $\zeta=1.0$。
\end{gongshi}

\textbf{突然缩小：}损失比突然扩大小得多。教材表 5.3（p83）按面积比给出 $\zeta$ 值
（$A_2/A_1=0.01\sim1.0$ 时 $\zeta=0.50\to0$，如 $0.1\to0.47$、$0.2\to0.45$、$0.5\to0.30$、$0.8\to0.15$）；
当 $A_2/A_1<0.10$ 时也可用经验式 $\zeta_2=0.5\left(1-\dfrac{A_2}{A_1}\right)$（$A_2$ 为小断面），
且 $\zeta$ 对应\emph{小断面}流速。

\begin{tishi}
\textbf{订正（以教材为准）：}原稿把 $\zeta_2=0.5\left(1-\dfrac{A_2}{A_1}\right)$ 说成突然缩小的通用经验式。
教材（表 5.3，p83）只在 $A_2/A_1<0.10$ 时用这条式子，一般面积比应按表 5.3 查取数值。
另外，突然扩大的损失公式教材称作\textbf{波达（Borda）公式}（p82），原稿写作「包达--卡诺公式」，
全章已统一为教材用语。
\end{tishi}

\begin{zhenti}{2015 年试题 计算第 3 题}{计算题}
水平突然缩小管路 $d_1=15\ \mathrm{cm}$、$d_2=10\ \mathrm{cm}$，水的流量 $Q=2\ \mathrm{m^3/min}$，
用水银测压计测得读数 $h=8\ \mathrm{cm}$，求突然缩小的水头损失。
\end{zhenti}
\noindent\textbf{思路：}① 由连续性方程求 $v_1$、$v_2$；② 用水银压差计读数换算两断面压差
$\dfrac{p_1-p_2}{\rho g}=\left(\dfrac{\rho_{Hg}}{\rho}-1\right)h$；③ 对两断面列伯努利方程，
水平管 $z_1=z_2$，得 $h_j=\dfrac{p_1-p_2}{\rho g}+\dfrac{v_1^2-v_2^2}{2g}$。

\textbf{常考题型：}计算（突然扩大/缩小的 $h_j$、$\zeta$）；填空（U 形水银压差计换算）；
选择（哪个管件 $\zeta$ 最大）。
\textbf{重要度 \xstarfull{4}}\quad\yiju{E4 2015 年试题计算第 3 题"突然缩小水头损失"；E4 2015 年试题填空"U 形水银压差计"；E1 slide33「第五章＝考试计算题重点」}\quad\nandu{中}

\begin{banfa}
\textbf{压差计换算通式（必背）：}$\dfrac{p_1-p_2}{\rho g}=\left(\dfrac{\rho'}{\rho}-1\right)h$，
$\rho'$ 为测压计工作液体（水银时 $\rho'/\rho=13.6$）。正接、倒接的区别只影响正负号。\\
\textbf{突然扩大做题三步：}① 算 $v_1$、$v_2$；② 判 $A_1$、$A_2$ 谁大；③ 直接代 $\dfrac{(v_1-v_2)^2}{2g}$，
不必求 $\zeta$。\\
\textbf{记住"扩大损失远大于缩小"}：这是出口损失 $\zeta=1.0$ 而进口只有 $0.5$ 的根本原因。
\end{banfa}
\end{kaodian}

\begin{kaodian}{3}{典型局部阻力系数与当量长度法}
\textbf{典型 $\zeta$ 值表（必记几个）：}
\begin{center}
\begin{tabularx}{\textwidth}{@{}lXX@{}}
\toprule
管件 & 局部阻力系数 $\zeta$（对应下游/小断面流速） & 说明 \\
\midrule
管道进口（锐缘） & 0.5 & 出口流入大容器时取 1.0 \\
管道出口（流入大容器） & 1.0 & 相当于突然扩大 $A_2\to\infty$ \\
突然扩大 & $\left(1-\dfrac{A_1}{A_2}\right)^2$ & 对应小断面流速 \\
突然缩小 & $0.5\left(1-\dfrac{A_2}{A_1}\right)$ & 对应小断面流速 \\
渐扩管 & 与扩张角、面积比有关 & 扩张角宜小于 $8^\circ$，角大则损失急增 \\
渐缩管 & 0.05$\sim$0.1（损失很小） & 收缩流稳定、不易脱离 \\
$90^\circ$ 弯头 & 0.5$\sim$1.0 & 与 $R/d$ 有关，$R/d$ 大则小 \\
闸阀（全开／半开） & 约 0.1$\sim$0.2／2$\sim$5 & 阀门开度是调节流量的主要手段 \\
\bottomrule
\end{tabularx}
\end{center}

\begin{gongshi}{当量长度法}
\[ h_j=\lambda\frac{l_e}{d}\frac{v^2}{2g},\qquad
   h_w=\lambda\frac{l+\sum l_e}{d}\frac{v^2}{2g} \]
\textbf{含义：}把局部损失折算成"同等损失的直管长度" $l_e$，称\textbf{当量长度}；
于是全管路损失统一写成沿程损失的形式，管件当作附加管长处理，算短管时非常方便。\\
\textbf{适用条件：}$l_e$ 与 $\lambda$ 有关的管件（如弯头、阀门）需按同一 $\lambda$ 折算；
进出口这类 $\zeta$ 与 $\lambda$ 无关的管件也可折算，但换算值随 $\lambda$ 变。
\end{gongshi}

\textbf{常考题型：}填空（"进口局部阻力系数为\rule{1.2cm}{0.4pt}"）；
计算（短管计算中把 $\sum\zeta$ 化为附加长度）。
\textbf{重要度 \xstarfull{3}}\quad\yiju{E1 slide33「第五章＝考试计算题重点」；E4 2015 年试题计算第 3 题（局部损失计算）；E2「第五章同为主要计算考点」}\quad\nandu{易}

\begin{banfa}
\textbf{两个"1 和 0.5"必背：}进口 $\zeta=0.5$、出口 $\zeta=1.0$——它们几乎出现在每道长管/短管的题里。
\textbf{"全开小、半开大"}：阀门开度越小 $\zeta$ 越大，节流调节就是靠它。
\textbf{当量长度法的价值：}$\dfrac{l+\sum l_e}{d}$ 一个分式搞定全管路，比逐项累加不易漏项。
\end{banfa}

\begin{yicuo}
① $\zeta$ 必须与它所乘的 $v$ \textbf{配套}：查表时先看清表里 $v$ 是"小断面流速"还是"出口流速"。
② 出口 $\zeta=1.0$ 不是"出口有阻力"，而是\emph{动能全部损失掉}，别理解成"出口被堵住"。
③ 渐扩管的损失通常\emph{大于}同面积比的突然扩大以外的一切管件中的渐缩管，
答题时不要因为"渐扩过渡平缓"就认为损失小。
\end{yicuo}
\end{kaodian}

\section{孔口、管嘴出流}

\begin{kaodian}{5}{薄壁小孔口自由出流}
\textbf{是什么：}器壁上开的孔洞叫\textbf{孔口}；液体经孔口泄出叫\textbf{孔口出流}。
孔口具有尖锐边缘、液体流过时只有线接触，称\textbf{薄壁孔口}，其阻力只有孔口处的局部阻力。
孔口高度 $e<\dfrac{1}{10}H$ 时为\textbf{小孔口}，可认为孔口断面上各点水头都等于孔口中心的水头；
出流于大气称\textbf{自由出流}，出流于液体中称\textbf{淹没出流}。

\noindent\textbf{教材原话（p84）：}「孔口出流与管嘴出流有一共同特点，即水流流出孔口或管嘴时局部损失起主导作用，
沿程损失可以略去不计。」教材的判别口径为：$\dfrac{d_0}{H}<0.1$ 时称\textbf{小孔口}，反之称大孔口；
孔口壁厚不大于孔径一半（$l\le d_0/2$）时称\textbf{薄壁孔口}；完善收缩的收缩系数 $\varepsilon=0.64$。

\begin{gongshi}{薄壁小孔口自由出流的流速与流量（教材式 5.38、5.39、5.40）}
\[ v_c=\varphi\sqrt{2gH_0},\qquad Q=A_cv_c=\varepsilon A\varphi\sqrt{2gH_0}=\mu A\sqrt{2gH_0},\qquad \mu=\varepsilon\varphi \]
\[ H_0=H+\frac{p_1-p_c}{\rho g}+\frac{v_1^2}{2g} \]
\textbf{含义：}$H_0$ 称\textbf{作用水头}；$\varphi=\dfrac{1}{\sqrt{1+\zeta_{\text{孔}}}}$ 称\textbf{流速系数}
（实际流速与理想流速之比）；$\varepsilon=\dfrac{A_c}{A}$ 称\textbf{收缩系数}（收缩断面面积与孔口面积之比）；
$\mu=\varepsilon\varphi$ 称\textbf{流量系数}（实际流量与理想流量之比）。
孔口的局部阻力系数 $\zeta_{\text{孔}}=\dfrac{1}{\varphi^2}-1$。\\
\textbf{适用条件：}薄壁、小孔口、自由出流、定水头（液面位置不变）、完善收缩。
自由出流时收缩断面压强为大气压，故 $p_1-p_c=0$，$H_0=H+\dfrac{v_1^2}{2g}$；
容器很大时 $v_1\approx0$，则 $H_0=H$。
\end{gongshi}

\textbf{薄壁圆形小孔口完善收缩的典型系数（必背）：}
$\varepsilon=0.63\sim0.64$，$\varphi=0.97\sim0.98$，$\mu=\varepsilon\varphi=0.60\sim0.62$。
取常见值 $\varepsilon\approx0.64$、$\varphi\approx0.97$、$\mu\approx0.62$。
教材（p85--86，教材式 5.40）给出的完善收缩取值即 $\varepsilon=0.64$、$\varphi=0.97$、$\mu=\varepsilon\varphi=0.62$，与上面取的常见值一致。
孔口 $\zeta_{\text{孔}}=\dfrac{1}{0.97^2}-1\approx0.06$，其他型式孔口的 $\zeta$ 都大于此值。

\textbf{常考题型：}计算（求 $Q$、$H$、孔口直径）；填空（"孔口流量系数的典型值"）；
选择（流速系数与流量系数的区别与大小关系）。
\textbf{重要度 \xstarfull{5}}\quad\yiju{E4 2026 回忆版计算第 5 题"管嘴体积流量（书上例题）"；E1 slide33「第五章＝考试计算题重点」；E2「第五章同为主要计算考点」}\quad\nandu{中}

\begin{banfa}
\textbf{一图记三个系数：}$\varepsilon$ 只管\emph{面积收缩}（$<1$），$\varphi$ 只管\emph{速度折减}（$<1$），
$\mu=\varepsilon\varphi$ 是二者的乘积，所以 $\mu<\varphi$ 且 $\mu<\varepsilon$ 之一不一定成立——
对孔口 $\mu\approx0.62$，比 $\varepsilon$ 略小、比 $\varphi$ 小很多。
\textbf{做题模板：}① 算 $H_0$；② 判孔口是大是小；③ 取 $\mu$ 代入 $Q=\mu A\sqrt{2gH_0}$；④ 反求 $H$ 时先平方。
\end{banfa}
\end{kaodian}

\begin{kaodian}{3}{大孔口出流与孔口淹没出流}
\textbf{大孔口：}孔口高度 $e\ge\dfrac{1}{10}H$ 时，孔口上下缘的水头差别不能忽略，
必须沿竖向积分求流量。设孔口上缘、下缘水头分别为 $H_1$、$H_2$，孔口宽 $b$，则
\[ Q=\frac{2}{3}\mu b\sqrt{2g}\left(H_2^{3/2}-H_1^{3/2}\right) \]
工程上常取孔口中心水头 $H_c=\dfrac{H_1+H_2}{2}$ 近似代入小孔口公式，误差不大。

\noindent\textbf{出处说明：}教材 5.6 节只给出薄壁小孔口自由出流（教材式 5.38$\sim$5.40）与孔口淹没出流
（教材式 5.41$\sim$5.43）；对大孔口只给出判别定义「当 $d_0/H<0.1$ 时，孔口称小孔口……反之则是大孔口」（教材 p84），
\emph{未列大孔口出流的积分式}。上面的大孔口公式取自外校辅导书口径，作为补充保留，考试范围以教材为准。

\textbf{孔口淹没出流（教材 5.6.2 节，教材式 5.41$\sim$5.43）：}孔口出流到下游液体中。教材先对 1-1、2-2 两断面
列伯努利方程（教材式 5.41），并说明淹没出流的流速系数「与自由出流的流速系数的表达式相同」，
从而得流速与流量式（教材式 5.42、5.43），形式上与小孔口自由出流完全相同：
\[ v_c=\varphi\sqrt{2gH},\qquad Q=\mu A\sqrt{2gH} \]
区别有二：① 此时 $H$ 是\textbf{上、下游液面差}（而不是孔口以上的水深）；
② 分母中的"1"代表突然扩大的局部阻力系数（出口动能损失在下游液体中）。
$\varepsilon$、$\varphi$、$\mu$ 的数值与自由出流相差不大，实用上认为相等。

\begin{gongshi}{变水头（非恒定）孔口出流（教材式 5.50）}
\[ T=\frac{2S\sqrt{H}}{\mu A\sqrt{2g}}=\frac{2SH}{\mu A\sqrt{2gH}} \]
\textbf{含义：}容器内液面不断下降、作用水头随时间变化的孔口出流称\textbf{变水头孔口出流}（教材 5.6.4 节）；
$S$ 为容器横断面面积、$A$ 为孔口面积、$H$ 为初始作用水头。教材指出该式「表明非恒定流的泄水时间
相当于相同水头作用下恒定泄放同体积水体所需时间的 2 倍」（教材 p88）——因为水头由 $H$ 降到零的过程中
流量由大变小，平均流量恰为恒定泄放流量的一半。\\
\textbf{适用条件：}容器断面 $S$ 在泄水过程中不变（或按平均断面处理）、孔口为小孔口且 $\mu$ 视为常数；
泄空至某一水头 $h$ 时，把式中的 $H$ 换成相应水头即可。
\end{gongshi}

\textbf{常考题型：}选择（"大孔口与小孔口的判别条件"$\to e<H/10$）；
填空（"淹没出流时 $H$ 代表\rule{2cm}{0.4pt}"）。
\textbf{重要度 \xstarfull{3}}\quad\yiju{E1 slide33「第五章＝考试计算题重点」；E2「第五章同为主要计算考点」（大孔口为次重点）}\quad\nandu{中}

\begin{yicuo}
① \textbf{自由出流与淹没出流最容易混}：自由出流的 $H$ 是\emph{孔口中心到上游液面的水深}；
淹没出流的 $H$ 是\emph{上下游液面高差}。写错这一点全题皆错。
② 大孔口的判别是 $e\ge H/10$，不是看孔口绝对尺寸（$10\ \mathrm{cm}$ 的孔在浅箱里也是大孔口）。
③ 淹没出流下游若不静止（有流速），还需在 $H$ 上加下游流速水头差。
\end{yicuo}
\end{kaodian}

\begin{kaodian}{5}{圆柱形外管嘴与正常工作条件}
\textbf{是什么：}在孔口上接一段短管（长度约为孔径的 $3\sim4$ 倍）就是\textbf{管嘴}。
液体进入管嘴后在距进口约 $0.5d$ 处形成\textbf{收缩断面}，随后又扩张充满整个管嘴再流出。

\noindent\textbf{教材原话（p87）：}「圆柱形外管嘴收缩断面的真空度可达到作用水头的 0.75 倍」；
「但当真空度达 7m 水柱以上时，液体会出现汽化现象，影响正常的管嘴出流」；「要使管嘴正常工作，
需同时满足以下两个条件：（1）作用水头 $H\le9\ \mathrm{m}$。（2）管嘴长度 $l=(3\sim4)d$。」

\begin{gongshi}{圆柱形外管嘴的出流公式与真空度（教材式 5.44~5.49）}
\[ v=\varphi\sqrt{2gH_0},\qquad Q=\mu A\sqrt{2gH_0},\qquad \mu=\varphi\approx0.82 \]
\[ \text{收缩断面真空度：}\ \frac{p_a-p_c}{\rho g}\approx0.75H_0 \]
\textbf{含义：}管嘴出口断面被液流充满，\textbf{没有收缩}，故 $\varepsilon=1$、$\mu=\varphi=0.82$，
比孔口的 $\mu\approx0.62$ 大得多（管嘴的"吸力"来自收缩断面的真空，相当于额外增加了作用水头）。
收缩断面处的真空度约为作用水头的 0.75 倍。\\
\textbf{适用条件：}\textbf{必须满足两个正常工作条件}：
① 作用水头 $H_0\le 9\ \mathrm{m}$（真空度不超过约 $7\ \mathrm{m}$ 水柱）；
② 管嘴长度 $l=(3\sim4)d$。两者\textbf{缺一不可}。
\end{gongshi}

\begin{jielun}
\textbf{为什么 $H_0\le9\ \mathrm{m}$：}由真空度 $\approx0.75H_0$，若 $H_0$ 太大，
收缩断面绝对压强降到液体汽化压强以下，液体汽化、外部空气被吸入管嘴，液流脱离管壁、
管嘴不再"充满"，$Q=\mu A\sqrt{2gH_0}$ 失效（$\mu$ 骤降到接近孔口值）。\\
\textbf{为什么 $l=(3\sim4)d$：}太短则液流尚未充满管嘴就已流出，收缩仍在，$\mu$ 接近孔口值；
太长则摩擦损失增大，性质变成短管。$3\sim4d$ 正是"刚好充满又摩擦不大"的长度。
\end{jielun}

\begin{zhenti}{2022 年 818 单选第 8 题}{单项选择题}
圆柱形直角外管嘴正常工作的条件是（\quad）。答：作用水头 $H_0\le9\ \mathrm{m}$，且管嘴长度 $l=(3\sim4)d$。
\end{zhenti}

\textbf{其他型式管嘴：}
\begin{center}
\begin{tabularx}{\textwidth}{@{}lXX@{}}
\toprule
型式 & 流量系数 $\mu$ & 特点与用途 \\
\midrule
圆柱形外管嘴 & 0.82 & 最常用，出口为满流，$Q$ 比孔口大 30\% 左右 \\
圆柱形内管嘴 & 约 0.71 & 管嘴伸入容器，水头损失大，$\mu$ 偏小 \\
圆锥形扩张管嘴 & 0.45$\sim$0.50 & $\mu$ 小但出口流速低、压强高，作喷射泵、水轮机尾水管 \\
圆锥形收缩管嘴 & 0.94$\sim$0.97 & 损失最小、$\mu$ 最大，用于消防水枪、水力采煤喷嘴 \\
流线形管嘴 & 约 0.98 & 损失最小，$\mu$ 最大，但加工复杂 \\
\bottomrule
\end{tabularx}
\end{center}

\textbf{孔口与管嘴流量系数之比（常考结论）：}$\dfrac{\mu_{\text{管嘴}}}{\mu_{\text{孔口}}}=\dfrac{0.82}{0.62}\approx1.32$，
即同样水头、同样断面积下管嘴比孔口多出约 32\% 的流量。
原因是孔口出口有收缩（$\varepsilon<1$）且收缩断面处无真空"增水头"；
管嘴出口无收缩（$\varepsilon=1$）且收缩断面存在约 $0.75H_0$ 的真空度，相当于增大了作用水头。

\begin{zhenti}{2026 年真题回忆版 计算第 5 题}{计算题}
求管嘴的体积流量（书上例题）。
\end{zhenti}
\noindent\textbf{思路：}先算 $H_0=H+\dfrac{v_1^2}{2g}$（容器大则 $H_0\approx H$），
再用 $Q=\mu A\sqrt{2gH_0}$，圆柱形外管嘴取 $\mu=0.82$。
\textbf{务必检查} $H_0\le9\ \mathrm{m}$ 与 $l=(3\sim4)d$；若 $H_0>9\ \mathrm{m}$，题目给的条件本身就不成立。

\textbf{常考题型：}选择（"圆柱形外管嘴正常工作条件"）；计算（管嘴流量、真空度）；
简答（孔口与管嘴流量系数为何不同）。
\textbf{重要度 \xstarfull{5}}\quad\yiju{E4 2022 年 818 单选第 8 题（圆柱形直角外管嘴正常工作条件 $H_0\le9\ \mathrm{m}$、$l=(3\sim4)d$）；E4 2026 回忆版计算第 5 题（管嘴体积流量，书上例题）}\quad\nandu{中}

\begin{banfa}
\textbf{"9 米、三四倍"口诀记正常条件}，并在答题时\emph{解释为什么}——这是拿满分的关键：
$H_0>9\ \mathrm{m}$ 时收缩断面真空度过大 $\to$ 空气吸入 $\to$ 管嘴失效。
\textbf{系数速查：}孔口 $\mu=0.62$、管嘴 $\mu=0.82$、收缩管嘴 $\mu=0.95$、流线形 $\mu=0.98$。
\textbf{"为什么管嘴流量大"的标准两句话：}① 出口无收缩，有效面积大；② 收缩断面真空形成"抽吸"作用，
相当于增大了作用水头。
\end{banfa}

\begin{yicuo}
① 管嘴正常工作的\textbf{两个条件缺一不可}，只写"$l=(3\sim4)d$"要丢分，只写"$H_0\le9\ \mathrm{m}$"也丢分，
必须两条并给出理由。
② $H_0=9\ \mathrm{m}$ 是\textbf{作用水头}，不是"孔口到液面的高度"——若容器有压，二者不同。
③ 管嘴 $\mu=\varphi=0.82$ 是因为 $\varepsilon=1$；孔口 $\mu=\varepsilon\varphi=0.62$ 是因为 $\varepsilon\approx0.64$。
④ 管嘴长度的单位是"$d$ 的倍数"，$l=(3\sim4)d$ 别误写成 $l=(3\sim4)$ m。
\end{yicuo}
\end{kaodian}

\section{管路的水力计算}

\begin{kaodian}{5}{长管与短管的判别}
\textbf{是什么：}把压力管路按计算中是否忽略局部损失与流速水头分成两类。
\begin{center}
\begin{tabularx}{\textwidth}{@{}lXX@{}}
\toprule
类别 & 定义（教材口径） & 典型场合 \\
\midrule
长管 & 与沿程水头损失相比，\textbf{流速水头和局部水头损失可以忽略}的管路 & 室外管路、油田地面集输管路、外输管路 \\
短管 & \textbf{流速水头和局部水头损失不能忽略}的管路 & 室内管路、联合站与计量间内管件较多的管路、水泵吸水管 \\
\bottomrule
\end{tabularx}
\end{center}

\noindent\textbf{教材原话（p88）：}「如管路系统的局部水头损失与速度水头的总和与沿程水头损失相比很小，
所占的比重通常小于 5\%，计算时可不考虑局部水头损失和速度水头，称为按长管计算」。教材并给出分类定义：
「等径无分支的管路系统称为简单管路。除简单管路外的管路系统为复杂管路。复杂管路通常按长管计算。」

\begin{tishi}
\textbf{订正（补教材的定量判据）：}教材该段末句排印作「否则为按长管计算」，按本段文意与 5.7 节内容，
第二个「长管」应为「短管」（这两句是同一判据的正反两面，教材此处系排印笔误）。原稿只写了
「长管和短管不按管道的绝对长度决定」这一定性说法，现补上教材的\textbf{定量判据}：
\textbf{局部水头损失与速度水头之和占沿程水头损失的比重小于 5\% 时按长管计算}，其余情况按短管计算。
\end{tishi}

\textbf{判别要点（必背）：}
① \textbf{长管和短管不按管道的绝对长度决定}——三里长的管子也可能是短管，三十米的也可能是长管；
② 管路上有较大局部损失管件（局部开启的闸门、喷嘴、底阀等）时，\emph{即使管道很长}局部损失也不能略去，必须按短管算；
③ 把管道按长管计算能简化计算过程，但在不能判断"流速水头与局部损失之和远小于沿程损失"之前，
按短管计算不会产生较大误差（即\emph{拿不准就当短管算}）。

\begin{gongshi}{长管的能量方程}
\[ z_1+\frac{p_1}{\rho g}=z_2+\frac{p_2}{\rho g}+h_f,\qquad
   H=\left(z_1+\frac{p_1}{\rho g}\right)-\left(z_2+\frac{p_2}{\rho g}\right)=h_f \]
\textbf{含义：}长管忽略流速水头与局部损失后，伯努利方程只剩位置水头、压强水头与沿程损失；
$H$ 称\textbf{作用水头}（管路两端的总水头之差），全部用于克服沿程阻力。\\
\textbf{适用条件：}长管（忽略 $\sum h_j$ 与 $\dfrac{v^2}{2g}$）；恒定流、液体充满全管、等直径段之间无分流。
因为忽略了流速水头，长管的\textbf{总水头线与测压管水头线重合}。
\end{gongshi}

\begin{zhenti}{2022 年 818 单选第 3 题}{单项选择题}
测压管水头线坡度小于零时，沿程各断面的测压管水头如何变化？
\end{zhenti}
\noindent\textbf{要点：}测压管水头线坡度 $=\dfrac{\dd}{\dd s}\left(z+\dfrac{p}{\rho g}\right)$，
小于零说明\emph{沿程下降}，反映机械能（对长管即总水头）沿程被阻力消耗。
\textbf{注意}：测压管水头线可以沿程上升（如渐扩管、突扩后），但\emph{总水头线一定沿程下降}。

\begin{zhenti}{题库 \texttt{10.pdf} 选择题第 9 题}{单项选择题}
有压管道的总水头线与测压管水头线的基本规律是（\quad）。\\
(A) 总水头线是沿程下降的。（B）测压管水头线沿程有可升可降。（C）测压管水头线是沿程下降的。（D）总水头线和测压管水头线沿程可升可降。
\end{zhenti}
\noindent\textbf{要点（答 A）：}总水头线因阻力耗能必须\textbf{沿程下降}；测压管水头线则可能升、可能降。

\textbf{常考题型：}选择/判断（长管短管判别、总水头线规律）；简答（"长管与短管的区别"）。
\textbf{重要度 \xstarfull{5}}\quad\yiju{E4 2022 年 818 单选第 3 题（测压管水头线坡度）；E5 题库选择第 9 题（总水头线沿程下降）；E5 题库计算第 2 题（水箱垂直管道出流 $Q$ 与管长 $l$ 的关系）}\quad\nandu{中}

\begin{banfa}
\textbf{判断口诀："长管只看 $h_f$，短管连 $v^2/2g$ 一起算。"}
\textbf{得分技巧：}答题时先写一句"本题按长管（或短管）计算"，这一步就是明确的给分点；
若题目没说明且管件很多，就按短管算并注明理由。
\textbf{总头线／测压管水头线的三条铁律：}总水头线\emph{只降不升}；测压管水头线\emph{可升可降}；
长管中二者\emph{重合}。
\end{banfa}
\end{kaodian}

\begin{kaodian}{5}{简单长管的三个基本问题}
\begin{gongshi}{长管的沿程损失与雷诺数（教材式 5.51）}
\[ h_f=\lambda\frac{l}{d}\frac{v^2}{2g}=0.0826\,\lambda\frac{Q^2l}{d^5},\qquad
   \Rey=\frac{4Q}{\pi d\nu},\qquad H=\frac{8\lambda LQ^2}{g\pi^2d^5} \]
\textbf{教材式（5.51）$H=\dfrac{8\lambda LQ^2}{g\pi^2d^5}$ 与上面的 $0.0826\,\lambda Q^2l/d^5$ 完全等价}
（$0.0826=8/(g\pi^2)$）。教材还引入\textbf{流量模数} $K$（教材式 5.52），把长管计算式简写成
$H=\dfrac{LQ^2}{K^2}$（教材式 5.53），$K$ 只与管径、$\lambda$ 有关，可直接查教材表 5.5
（铸铁管的流量模数 $K$ 值，p89）。
\textbf{含义：}把 $v=\dfrac{4Q}{\pi d^2}$ 代入达西公式即得用流量表示的形式；
雷诺数直接用流量表达 $\Rey=\dfrac{4Q}{\pi d\nu}$。层流时 $\lambda=\dfrac{64}{\Rey}=\dfrac{16\pi d\nu}{Q}$，
代回得 $h_f\approx4.15\,\dfrac{\nu Ql}{d^4}$（层流长管专用，与 $d^4$ 成反比）。\\
\textbf{适用条件：}圆管、恒定流；层流式仅当 $\Rey<2300$；湍流时 $\lambda$ 由第 5.4 节分区公式确定。
\end{gongshi}

\textbf{三个基本问题：}
\begin{center}
\begin{tabularx}{\textwidth}{@{}lXX@{}}
\toprule
问题 & 已知 & 求解方法与难点 \\
\midrule
第一类 & $d$、$l$、$Q$（及流体性质） & 求 $h_f$。可直接计算：算 $v$、$\Rey$、判区、取 $\lambda$、代入达西式 \\
第二类 & $d$、$l$、$h_f$ & 求 $Q$。$\Rey$ 含 $Q$、$\lambda$ 又含 $\Rey$，须用\textbf{试差法或迭代} \\
第三类 & $l$、$Q$、$h_f$ & 求 $d$。$d$ 隐含在 $\Rey$ 与 $\Delta/d$ 中，同样须试差；工程上先选定\textbf{经济流速}再选标准管径 \\
\bottomrule
\end{tabularx}
\end{center}

\textbf{经济流速：}流速过大，管子易磨损、迅速关阀时易产生较大水击压力而破裂；
流速过小，液体中杂质易在管中沉淀、侵蚀管壁并增大阻力。一般油田内部管路取 $1\sim2\ \mathrm{m/s}$，
外输管路取 $1\sim3\ \mathrm{m/s}$ 为宜。

\begin{zhenti}{题库 \texttt{10.pdf} 计算第 2 题}{计算题}
水箱中的水通过垂直管道向大气出流。设水箱水深为 $H$，管道直径 $d$，长度 $l$，沿程阻力系数 $\lambda$，
局部阻力系数 $\zeta$。试求：① 在什么条件下流量 $Q$ 不随管长 $l$ 而变？② 什么条件下 $Q$ 随管长 $l$ 的加大而增加？
③ 什么条件下 $Q$ 随管长 $l$ 的加大而减小？
\end{zhenti}
\noindent\textbf{思路（按短管列能量方程）：}由液面到出口断面列伯努利方程，
$H+l=\left(\lambda\dfrac{l}{d}+\zeta+1\right)\dfrac{v^2}{2g}$，得
$v=\sqrt{\dfrac{2g(H+l)}{\lambda l/d+\zeta+1}}$。
讨论 $\dfrac{\partial Q}{\partial l}$ 的符号即得三种情形：
① 当 $H=0$（或 $H$ 与 $l$ 的组合使分子分母同比例变化）时 $Q$ 与 $l$ 无关；
② 当 $H$ 相对 $l$ 起主导（$\zeta$ 很大、$\lambda l/d$ 很小）时 $Q$ 随 $l$ 增大而增大；
③ 当 $l$ 起主导时 $Q$ 随 $l$ 增大而减小。
\textbf{这题的关键是：把"水箱水深 $H$"与"管长 $l$"都放进作用水头，并用 $\zeta$ 与 $\lambda l/d$ 的相对大小讨论。}

\textbf{常考题型：}计算（三类问题之一，第二、三类多需迭代）；
简答（"为什么第二、三类问题要用试差法"）。
\textbf{重要度 \xstarfull{5}}\quad\yiju{E5 题库计算第 2 题（水箱垂直管道出流 $Q$ 与管长 $l$ 的关系）；E1 slide33「第五章＝考试计算题重点」；E2「第五章同为主要计算考点」}\quad\nandu{难}

\begin{banfa}
\textbf{第一类问题 5 步模板：}① $v=\dfrac{4Q}{\pi d^2}$；② $\Rey=\dfrac{vd}{\nu}$；③ 判别层流／湍流及分区；
④ 取 $\lambda$；⑤ $h_f=\lambda\dfrac{l}{d}\dfrac{v^2}{2g}$。\\
\textbf{第二、三类问题的迭代写法（拿过程分）：}设 $\lambda_0=0.03\to$ 求 $Q_0\to$ 算 $\Rey_0\to$ 查新 $\lambda_1\to$ 求 $Q_1$，
一般两次即收敛。答题时把两次结果都写上，并注明"已收敛"。\\
\textbf{单位速查：}$1\ \mathrm{L/s}=10^{-3}\ \mathrm{m^3/s}$；$\nu$ 给 $\mathrm{mm^2/s}$ 要换成 $\mathrm{m^2/s}$（除以 $10^6$）。
\end{banfa}
\end{kaodian}

\begin{kaodian}{5}{串联管路与并联管路}
\begin{gongshi}{串联管路的水力特征}
\[ Q_1=Q_2=\cdots=Q_n=Q,\qquad h_f=\sum_{i=1}^{n}h_{fi}=h_{f1}+h_{f2}+\cdots+h_{fn} \]
\textbf{含义：}串联管路各段流量\textbf{相同}（连续性原理）；全线水头损失为各段\textbf{之和}（能量守恒原理）。
节点处满足 $\sum Q_i=0$（进为正、出为负）。\\
\textbf{适用条件：}恒定流、管段之间无分流（无节点流出）；各段可分别用各自的 $d_i$、$l_i$、$\lambda_i$ 算 $h_{fi}$。
\end{gongshi}

\begin{gongshi}{并联管路的水力特征}
\[ Q=\sum_{i=1}^{n}Q_i=Q_1+Q_2+\cdots+Q_n,\qquad
   h_{f1}=h_{f2}=\cdots=h_{fn}=h_{A\to B} \]
\textbf{含义：}并联管路各支管的\textbf{水头损失相等}（都是同一对节点 $A$、$B$ 之间的能量差），
总流量为各支管\textbf{流量之和}。\\
\textbf{适用条件：}恒定流；各支管的两端节点相同；不计节点处的局部损失或已把它计入支管。
\end{gongshi}

\textbf{串并联求解套路：}
\begin{enumerate}
\item 串联：从总 $Q$ 出发，逐段算 $v_i$、$\Rey_i$、$\lambda_i$、$h_{fi}$，最后累加即得所需水头 $H$。
\item 并联：先写 $h_{f1}=h_{f2}$ 的比例式。由 $h_f=0.0826\lambda\dfrac{Q^2l}{d^5}$ 得
$\dfrac{Q_1}{Q_2}=\sqrt{\dfrac{\lambda_2l_2d_1^5}{\lambda_1l_1d_2^5}}$，求出流量分配比，再配合 $Q_1+Q_2=Q$ 解出各支管流量，
最后代回任一式求 $h_{A\to B}$。
\end{enumerate}

\textbf{串、并联在长输管线上的用途：}① 增加输送流量（并联）；② 延伸输送距离（串联）；
③ 克服翻越点。

\textbf{常考题型：}计算（并联流量分配、串联总水头）——\textbf{本章最高频的大题}；
选择（"并联管路各支管的水头损失相等还是流量相等"）。
\textbf{重要度 \xstarfull{5}}\quad\yiju{E1 slide33「第五章＝考试计算题重点」；E2「第五章同为主要计算考点，需掌握计算套路」；E4 2026 回忆版计算第 5 题（管路类计算）}\quad\nandu{难}

\begin{banfa}
\textbf{一句话记住串并联：串联"流量相等、损失相加"；并联"损失相等、流量相加"。}
反过来记是最常见的失分点，考场上先写下这句话再列式。\\
\textbf{并联题必用比例式：}$Q_i\propto\sqrt{\dfrac{d_i^5}{\lambda_i l_i}}$（$\lambda$ 相同时 $\propto\sqrt{d^5/l}$）。
由比例式 + 总流量两个方程即可解出全部 $Q_i$，比设未知数迭代快得多。\\
\textbf{并联题的检查：}$\sum Q_i$ 必须等于总流量；各支管算出的 $h_f$ 必须相同——用这两条自查。
\end{banfa}

\begin{yicuo}
① \textbf{并联读成"流量相等"、串联读成"损失相等"}是历年最高频的概念失分点。
② 并联支管中\textbf{短管阻力小的那条流量大}，不是"管径大流量就大"（还与 $l$、$\lambda$ 有关）。
③ 串联管路中若某段突然扩大，该处还有局部损失，别只累加沿程损失。
④ 并联管路的两支管长度、直径不同也可以，唯一必须相同的是\emph{两端节点}。
\end{yicuo}
\end{kaodian}

\begin{kaodian}{4}{分支管路与管网}
\textbf{是什么：}自某一点分开后\emph{不再汇合}的管路叫\textbf{分支管路}；由多条支管、环状管段组成的系统叫\textbf{管网}。

\textbf{分支管路的水力特征：}① 节点处流量平衡（流入＝流出）；
② 沿任一条干线，总水头损失＝该线各段损失之和 $h_f=\sum h_{fi}$；
③ 节点处 $\left(z+\dfrac{p}{\gamma}\right)$ 为定值（测压管水头在节点唯一）。

\textbf{分支管路的求解特点：}与并联的"损失相等"不同，分支管路各支管\emph{流量未知且互不相等}，
通常以\textbf{节点为枢纽}：先求各支管到节点所需的测压管水头，再由节点水头与各支管阻力反求流量，
并用节点流量平衡校验。

\textbf{常考题型：}填空（分支管路节点流量平衡）；计算（给节点水头求各支管流量）。
\textbf{重要度 \xstarfull{4}}\quad\yiju{E1 slide33「第五章＝考试计算题重点」；E2「第五章同为主要计算考点」}\quad\nandu{难}

\begin{banfa}
\textbf{节点法三步：}① 在节点列 $\sum Q=0$；② 由已知支管反求节点测压管水头
$z+\dfrac{p}{\gamma}$；③ 用该水头算其余支管流量。\\
\textbf{与并联的区别一句话：}并联是\emph{两端都固定}（损失相等），分支是\emph{只有一端固定}（节点水头相等）。
\end{banfa}
\end{kaodian}

\begin{kaodian}{4}{短管的水力计算、虹吸管与水泵管路}
\begin{gongshi}{短管的实用计算通式}
\[ H_0=z_1-z_2+\frac{p_1-p_2}{\rho g}+\frac{v_1^2}{2g}=(1+\zeta_c)\frac{v_2^2}{2g} \]
\[ Q=\mu A\sqrt{2gH_0},\qquad \mu=\frac{1}{\sqrt{1+\zeta_c}},\qquad
   \zeta_c=\sum\lambda_i\frac{l_i}{d_i}\left(\frac{v_i}{v_2}\right)^2+\sum\zeta_j\left(\frac{v_j}{v_2}\right)^2 \]
\textbf{含义：}短管必须把沿程损失、局部损失和出口流速水头全部计入。
$\zeta_c$ 称\textbf{综合阻力系数}，做法是\emph{一切损失都折算成出口流速 $v_2$ 的动能倍数}：
不同管径段的 $v_i$ 用连续性方程 $v_i=v_2\left(\dfrac{d_2}{d_i}\right)^2$ 折算。
$\mu$ 称短管的流量系数。\\
\textbf{适用条件：}短管（不能忽略 $\sum h_j$ 与 $\dfrac{v^2}{2g}$）；
自由出流时 $H_0=H+\dfrac{v_1^2}{2g}$（$H$ 为上游液面到出口中心的高差），
淹没出流时 $H_0$ 为\textbf{上下游液面差}。
\end{gongshi}

\textbf{水泵吸水管与汽蚀：}列自吸水池液面到水泵进口断面的伯努利方程，可得安装高度
\[ [h_s]=\frac{p_a}{\rho g}-\frac{p_v}{\rho g}-h_w-\frac{v^2}{2g} \]
（$p_v$ 为液体在当时温度下的汽化压强）。水泵进口压强必须高于 $p_v$，否则发生\textbf{汽蚀}——
这与泵的\textbf{允许吸上真空高度 $H_s$} 是同一件事的两种说法。

\textbf{虹吸管：}外形呈倒 U 形、靠上下游液面高差驱动的管路。计算两步：
① 对上下游液面列能量方程求流量 $Q$ 与管内流速；
② 对\emph{最高断面}列能量方程求该处压强，要求\textbf{最高处压强不低于液体的汽化压强}——
即最高处出现\emph{负压（真空）}，通常限制其真空度不超过 $7\sim8\ \mathrm{m}$ 水柱。
所以虹吸管的最高点不能离上游液面太高。

\begin{zhenti}{2022 年 818 单选第 7 题、题库 \texttt{10.pdf} 选择第 8 题}{单项选择题}
虹吸管最高处的压强（\quad）。（A）大于大气压（B）等于大气压（C）小于大气压（D）无法确定
\end{zhenti}
\noindent\textbf{要点（答 C）：}虹吸管最高处位于上游液面以上，能量方程给出
$\dfrac{p}{\rho g}=\dfrac{p_a}{\rho g}-h-\dfrac{v^2}{2g}-h_w<0$（相对压强为负），故\textbf{小于大气压}。
这一条也是"最高点不能太高"的原因。

\begin{zhenti}{题库 \texttt{10.pdf} 计算第 4 题}{计算题}
一台水泵 $n=1450\ \mathrm{r/min}$，$H$--$Q$ 性能曲线由表给出；管路系统综合阻力数 $S=0.024\times10^6\ \mathrm{s^2/m^5}$，
几何扬水高度 $H_z=6\ \mathrm{m}$，上下两水池水面均为大气压。求：（1）水泵运行的工作参数 $(Q,H,N,\eta)$；
（2）可采取哪些方式调节运行工况？（3）节流阀调节使 $Q=8\ \mathrm{L/s}$ 时各工作参数。
\end{zhenti}
\noindent\textbf{要点：}管路特性曲线 $H=H_z+SQ^2=6+0.024\times10^6Q^2$；
与泵的 $H$--$Q$ 曲线的\textbf{交点即工况点}，读出 $Q$、$H$，再由 $\eta$ 算轴功率。
调节方式：节流调节、变速调节、切割叶轮、并联或串联泵。节流后工况点沿泵曲线左移，$H$ 增大、$Q$ 减小、效率下降。

\textbf{常考题型：}计算（短管流量、吸水管安装高度、虹吸管最高处压强、水泵工况点）；
选择（虹吸管最高处压强、允许吸上真空高度含义）。
\textbf{重要度 \xstarfull{4}}\quad\yiju{E4 2022 年 818 单选第 7 题（虹吸管最高处压强）；E5 题库选择第 8 题（虹吸管最高处压强小于大气压）、第 14 题（允许吸上真空高度 $H_s$）；E5 题库计算第 4 题（水泵工况点与节流调节）}\quad\nandu{难}

\begin{banfa}
\textbf{短管计算三步：}① 写能量方程，左右各写全（含 $\dfrac{v^2}{2g}$ 与 $\sum h_j$）；
② 把所有速度折算成出口速度，凑成 $\zeta_c$；③ 用 $Q=\mu A\sqrt{2gH_0}$ 求流量。\\
\textbf{虹吸管两问两答：}问流量$\to$用上下游液面高差；问最高处压强$\to$以上游液面到最高点的高差列式。
两个方程不要用混。\\
\textbf{汽蚀一句话：}进口压强降到汽化压强以下 $\to$ 液体汽化形成气泡 $\to$ 气泡在高压区溃灭冲击叶片 $\to$ 汽蚀破坏。
\end{banfa}
\end{kaodian}

\begin{liti}{5.3}{书 p88--89 · 短管·水泵管路·安装高度与提水高度}
\textbf{题目}（教材原题，p88）水泵将水自池抽至水塔，如图 5.22 所示。已知：水泵的功率 $P=25\ \mathrm{kW}$，
流量 $Q=0.06\ \mathrm{m^3/s}$，水泵效率 $\eta=75\%$；吸水管长度 $l_1=8\ \mathrm{m}$，压水管长度 $l_2=50\ \mathrm{m}$，
吸水管直径 $d_1=250\ \mathrm{mm}$，压水管直径 $d_2=200\ \mathrm{mm}$，沿程阻力系数 $\lambda=0.025$；
带底阀滤水网的局部阻力系数 $\zeta=4.4$，弯头阻力系数 $\zeta=0.2$（2 个），阀门 $\zeta=0.5$，
单向阀 $\zeta=5.5$；水泵的允许真空度 $h_v=6\ \mathrm{m}$。试求：（1）水泵的安装高度；（2）水泵的提水高度。

\tuyi{吸水池内的水经吸水管（$d_1=250\ \mathrm{mm}$、$l_1=8\ \mathrm{m}$，管口带底阀滤水网）被水泵吸入，
再经压水管（$d_2=200\ \mathrm{mm}$、$l_2=50\ \mathrm{m}$，管路上有弯头、阀门与逆止阀）送入水塔。
水泵轴线高出吸水池水面的竖直距离即安装高度 $h_s$；水塔水面与吸水池水面的高差即提水高度。}

\begin{jie}
\textbf{① 求两段管中的流速。}由连续性方程 $v=\dfrac{4Q}{\pi d^2}$，吸水管中
\[ v_1=\frac{4Q}{\pi d_1^2}=\frac{4\times0.06}{\pi\times0.25^2}=1.22\ \mathrm{m/s} \]
压水管中
\[ v_2=\frac{4Q}{\pi d_2^2}=\frac{4\times0.06}{\pi\times0.2^2}=1.91\ \mathrm{m/s} \]
对应的流速水头为 $v_1^2/2g=0.0759\ \mathrm{m}$、$v_2^2/2g=0.186\ \mathrm{m}$（取 $g=9.8\ \mathrm{m/s^2}$）。\\
\textbf{② 列吸水池液面 0--0 到水泵进口断面 1--1 的伯努利方程。}以吸水池液面为基准面，
液面面积很大故 $v_0\approx0$、其压强为大气压，水泵进口是吸水管上压强最低的断面，得
\[ 0+\frac{p_a}{\rho g}+0=h_s+\frac{p_1}{\rho g}+\frac{v_1^2}{2g}+h_{w0\to1} \]
即
\[ h_s=\frac{p_a-p_1}{\rho g}-\frac{v_1^2}{2g}-h_{w0\to1} \]
式中 $\dfrac{p_a-p_1}{\rho g}$ 就是\textbf{允许真空度} $h_v=6\ \mathrm{m}$，故
\[ h_s=h_v-\frac{v_1^2}{2g}-h_{w0\to1} \]
\textbf{③ 计算吸水管的水头损失。}吸水管上有沿程损失与两项局部损失（带底阀滤水网 $\zeta=4.4$、
弯头 $\zeta=0.2$），全部按 $v_1$ 的流速水头折算：
\[
\begin{aligned}
h_{w0\to1}&=\left(\lambda\frac{l_1}{d_1}+\sum\zeta\right)\frac{v_1^2}{2g}
   =\left(0.025\times\frac{8}{0.25}+4.4+0.2\right)\times0.0759\\
   &=(0.8+4.6)\times0.0759=0.41\ \mathrm{m}
\end{aligned}
\]
代回即得\textbf{安装高度}
\[ h_s=6-0.0759-0.41=5.51\ \mathrm{m} \]
水泵安装高度不得超过 $5.51\ \mathrm{m}$，否则泵进口压强降到汽化压强以下，会发生\textbf{汽蚀}。\\
\textbf{④ 求水泵扬程。}水泵的有效功率 $P_u=\rho gQH$ 与轴功率 $P$ 的关系为 $P_u=\eta P$，故
\[ H=\frac{\eta P}{\rho gQ}=\frac{0.75\times25\times10^3}{1000\times9.8\times0.06}=31.89\ \mathrm{m} \]
\textbf{⑤ 计算压水管的水头损失。}压水管上除沿程损失外，还有弯头 $\zeta=0.2$、阀门 $\zeta=0.5$、
单向阀 $\zeta=5.5$，以及\textbf{出口损失} $\zeta=1.0$（出口动能全部损失在水塔里），
全部按 $v_2$ 的流速水头折算：
\[
\begin{aligned}
h_{w1\to2}&=\left(\lambda\frac{l_2}{d_2}+0.2+0.5+5.5+1.0\right)\frac{v_2^2}{2g}\\
   &=(6.25+7.2)\times0.186=2.50\ \mathrm{m}
\end{aligned}
\]
于是从吸水池液面到水塔液面（0--0 到 2--2）的全部水头损失为
\[ h_{w0\to2}=h_{w0\to1}+h_{w1\to2}=0.41+2.50=2.91\ \mathrm{m} \]
\textbf{⑥ 求提水高度。}列 0--0、2--2（水塔液面）断面的伯努利方程，两断面的压强水头都是大气压、
速度水头都近似为零，$H$ 为水泵扬程，设提水高度为 $H_z$，则
\[ 0+0+0+H=H_z+0+0+h_{w0\to2} \]
\[ H_z=H-h_{w0\to2}=31.89-2.91=28.98\ \mathrm{m} \]
即\textbf{提水高度为 28.98 m}。可见水泵扬程 $31.89\ \mathrm{m}$ 中约 $28.98\ \mathrm{m}$ 用于提升水体，
其余 $2.91\ \mathrm{m}$ 用于克服全管路的水头损失。
\end{jie}

\begin{tishi}
\textbf{OCR 不可判读与订正（本刊处理）：}OCR（\texttt{02-ocr/zhqin/page-0097.txt}、\texttt{page-0098.txt}）
把分式、下标、希腊字母整片丢失，题干印作「功率 P。=25kW」「效率 m=75%」「流量 Q=0.06m²/s」
「沿程阻力系数入=0.025」「允许真空度 h.=6m」，解答只剩「=6-6.4」「=0.41」「h=5.51m」
「PM/(pgQ)=31.89m」「=0.41+2.5=2.91m」「h=31.89-2.91=28.98m」这类碎片。
本讲义按\textbf{伯努利方程 + 短管综合阻力系数}反推重建，并用三个锚点校核：
（i）安装高度 $h_s=6-\dfrac{v_1^2}{2g}-0.41=5.51\ \mathrm{m}$ 可复现；
（ii）扬程 $H=\dfrac{\eta P}{\rho gQ}=31.89\ \mathrm{m}$ 可复现；
（iii）$0.41+2.50=2.91\ \mathrm{m}$ 可复现。
由此可确定两处教材未明写的事项：\textbf{① 2 个弯头按吸水管、压水管各 1 个分配}
（只有这样才能同时得到 0.41 m 与 2.50 m 两个损失值）；\textbf{② 压水管的出口损失取 $\zeta=1.0$}
（即出口流速水头 $v_2^2/2g$ 全部计入损失，故系数和为 $6.25+7.2$ 而不是 $6.25+6.2$）。
带底阀滤水网的 $\zeta=4.4$ 在扫描件上偏模糊，\textbf{照教材数值保留、未改}；
$\eta$（原印作「m」）、$\lambda$（原印作「入」）属符号订正，\textbf{数值一律未改}；
功率下标「P。」不可判读，本讲义按题干数值记作轴功率 $P=25\ \mathrm{kW}$。\\
\textbf{流速订正：}本题两段流速极易写反——吸水管管径\textbf{大}（$d_1=250\ \mathrm{mm}$）故流速\textbf{小}，
$v_1=1.22\ \mathrm{m/s}$；压水管管径\textbf{小}（$d_2=200\ \mathrm{mm}$）故流速\textbf{大}，$v_2=1.91\ \mathrm{m/s}$。
本讲义\textbf{不采信}把二者对调成 $v_1=1.91\ \mathrm{m/s}$、$v_2=1.22\ \mathrm{m/s}$ 的写法；
两个结果数字（安装高度 $5.51\ \mathrm{m}$、提水高度 $28.98\ \mathrm{m}$）与教材一致。
\end{tishi}
\end{liti}

\section{管路中的水击}

\begin{kaodian}{4}{水击现象、相长与直接／间接水击}
\textbf{是什么：}由于某种原因引起管内液体流速\emph{突然变化}（迅速开关阀门、突然停泵等），
管内压强随之急剧升降，并以压力波的形式在管内往复传播，这种现象叫\textbf{水击现象}，也叫\textbf{水锤}。
阀门突然关闭时流速骤减、压强首先急剧升高，称\textbf{正水击}；阀门突然开启时流速骤增、压强首先降低，称\textbf{负水击}。

\noindent\textbf{教材原话（p91）：}「在有压管路流动中，由于阀门突然开启或关闭等原因，导致流速发生急剧变化，
引起压强的剧烈波动，并在整个管长范围内传播的现象，称为水击或水锤。」教材并指出水击的危害：
「会导致管道强烈的震动、噪声和气穴，有时甚至引起管道变形、爆裂或阀门的破坏」。

\textbf{发生水击的物理原因（简答要点）：}
① \textbf{外因}——管路中流速突然变化；
② \textbf{内因}——液体具有\textbf{惯性}和\textbf{压缩性}，管壁具有\textbf{弹性}。
惯性力图维持液体原来的运动状态，压缩性与管壁弹性则改变体积、缓和流动。
水击中的液流参数随时间和位置变化，属于\textbf{不稳定流动}。

\begin{gongshi}{水击的相长、直接水击与水击压强（教材式 5.54、5.56、5.57）}
\[ T=\frac{2L}{c},\qquad
   t\le T\ \Rightarrow\ \text{直接水击};\quad t>T\ \Rightarrow\ \text{间接水击} \]
\[ \Delta p=\rho c\left(v_0-v\right) \]
\textbf{含义：}$L$ 为管道长度、$c$ 为水击波传播速度，$T=\dfrac{2L}{c}$ 称\textbf{相长}（压力波往返一次的时间）。
\textbf{直接水击：}阀门关闭时间 $t\le T$，压力波尚未返回，阀门前达到最大升压；
\emph{直接水击危害最大}。\textbf{间接水击：}关闭时间 $t>T$，压力波在关闭过程中已部分返回叠加，
$\Delta p$ 小于直接水击，故\textbf{延长关闭时间就能把直接水击变成间接水击}。
水击压强增量由儒科夫斯基公式 $\Delta p=\rho c(v_0-v)$ 给出，$v$ 为关闭后的流速（完全关闭时 $v=0$，
$\Delta p=\rho c v_0$）。\\
\textbf{适用条件：}直接水击公式只适用于 $t\le T$；考虑液体压缩性与管壁弹性；
$\Delta p$ 是压强\emph{增量}，实际压强为 $p_0+\Delta p$。
\end{gongshi}

\noindent\textbf{教材原话（p93）：}「水击波从阀门传播到管道进口再返回阀门所需时间称为水击的相，
以 $t_r$ 表示，两相为一个周期」；阀门关闭时间 $t\le t_r$ 为\textbf{直接水击}，$t>t_r$ 为\textbf{间接水击}。

\textbf{水击物理过程的四个传播阶段（教材表 5.6，p93）：}
\begin{center}
\begin{tabularx}{\textwidth}{@{}lXX@{}}
\toprule
阶段 & 时间 & 流速、压强与管壁的变化 \\
\midrule
I 增压逆波 & $0<t<\dfrac{L}{c}$ & 流速 $v_0\to0$，压强 $p_0\to p_0+\Delta p$，管壁膨胀 \\
II 降压顺波 & $\dfrac{L}{c}<t<\dfrac{2L}{c}$ & 流速 $0\to-v_0$，压强 $p_0+\Delta p\to p_0$，管壁恢复原状 \\
III 降压逆波 & $\dfrac{2L}{c}<t<\dfrac{3L}{c}$ & 流速 $-v_0\to0$，压强 $p_0\to p_0-\Delta p$，管壁收缩 \\
IV 增压顺波 & $\dfrac{3L}{c}<t\le\dfrac{4L}{c}$ & 流速 $0\to v_0$，压强 $p_0-\Delta p\to p_0$，管壁恢复原状 \\
\bottomrule
\end{tabularx}
\end{center}

\textbf{常考题型：}简答（"什么是水击现象？原因与危害？"）——2026 回忆版原题；
选择/填空（相长公式、直接与间接水击的判别、水击压强的增减方向）。
\textbf{重要度 \xstarfull{4}}\quad\yiju{E4 2026 回忆版简答考"水击现象"；E1 slide33「第五章＝考试计算题重点」；E2「第五章同为主要计算考点」}\quad\nandu{中}

\begin{banfa}
\textbf{简化式：}$t\le\dfrac{2L}{c}$ 直接水击（$\Delta p$ 最大），$t>\dfrac{2L}{c}$ 间接水击。
\textbf{简答答题模板（四句）：}① 定义（流速突变$\to$压强急剧升降$\to$波在管内传播）；
② 原因（外因流速突变，内因液体的惯性＋压缩性＋管壁弹性）；③ 分类（正／负水击，直接／间接水击）；
④ 危害与防护。\\
\textbf{正负水击速判：}关阀门$\to$流速减小$\to$压强升高（正水击）；开阀门$\to$流速增大$\to$压强降低（负水击）。
\end{banfa}

\begin{yicuo}
① 湍流不是水击的前提，水击处理的是\textbf{不稳定流动}，与第 5.2 节的层流湍流是两套体系。
② $\Delta p=\rho c(v_0-v)$ 中 $c$ 是\textbf{压力波传播速度}（约 $1000\ \mathrm{m/s}$ 量级），
不是管内流速，两者相差三个数量级。
③ 相长是"往返一次"的时间，故为 $\dfrac{2L}{c}$，不是 $\dfrac{L}{c}$。
④ 直接水击才用 $\Delta p=\rho c v_0$；间接水击的 $\Delta p$ 必须由具体关闭过程求出，不能直接套。
\end{yicuo}
\end{kaodian}

\begin{kaodian}{4}{水击波速、危害与防护措施}
\begin{gongshi}{水击波传播速度（教材式 5.55）}
\[ c=\frac{\sqrt{K/\rho}}{\sqrt{1+\dfrac{DK}{\delta E}}} \]
\textbf{含义：}$\sqrt{K/\rho}$ 是液体中的声速（$K$ 为液体体积弹性模量）；
分母中的 $\dfrac{DK}{\delta E}$ 是对\textbf{管壁弹性}的修正（$D$ 为管径、$\delta$ 为管壁厚度、$E$ 为管材弹性模量）。
管壁越硬、管径越小、壁越厚，$c$ 越接近声速；
管壁有弹性会吸收一部分能量，使 $c$ 减小、水击压强减小（这正是采用弹性管道或加设弹性元件能减轻水击的机理）。
\\
\textbf{适用条件：}薄壁圆管、管壁弹性与液体压缩性同时考虑；若管壁视为绝对刚性，则 $c=\sqrt{K/\rho}$。
\end{gongshi}

\textbf{教材算例数据（p93）：}$K=20.6\times10^9\ \mathrm{N/m^2}$；管材弹性模量 $E$：钢管 $19.6\times10^{10}$、
铸铁管 $9.8\times10^{10}$、混凝土管 $20.58\times10^9\ \mathrm{N/m^2}$。在 $d=0.2\ \mathrm{m}$、$\delta=0.01\ \mathrm{m}$
的钢管中 $c=1305\ \mathrm{m/s}$，教材据此指出「水击波的传播速度是很高的」。

\textbf{间接水击的近似计算（教材式 5.58，p94）：}$\Delta p\approx\rho c v_0\dfrac{t_r}{t_s}$
（$t_s$ 为阀门关闭时间）。教材指出「由上式可以看出 $t_r$ 越大，则 $\Delta p$ 越小」——
这就是延长关闭时间能减小水击压强的定量依据。教材列出的减小水击危害措施有三条（p94）：
延长阀门关闭时间；减小管中的流速；设置安全阀、水击消除阀等。

\textbf{水击的危害与防护措施（简答必背）：}
\begin{center}
\begin{tabularx}{\textwidth}{@{}lXX@{}}
\toprule
环节 & 具体做法 & 原理 \\
\midrule
延长关闭时间 & 把 $t$ 拉长到 $t>\dfrac{2L}{c}$，使直接水击变为间接水击 & 压力波在关闭过程中返回叠加，抵消部分升压 \\
设空气阀／空气包 & 在阀门或管道前装空气包、空气阀 & 气体可压缩，吸收水击能量 \\
设调压塔 & 在管道中间设调压塔（井） & 提供自由液面，压力波在此反射、泄压 \\
装水锤消除器 & 在阀门前装安全阀、水锤消除器 & 超过允许压强时泄水，保护管路 \\
加大管径 & 降低管内流速 $v_0$ & $\Delta p=\rho c(v_0-v)\propto v_0$，流速低则升压小 \\
缩短管路 & 减小 $L$，即减小相长 $T$ & 缩短压力波往返时间 \\
\bottomrule
\end{tabularx}
\end{center}

\textbf{常考题型：}简答（"水击的危害与防护措施"）；填空（水击波速公式的影响因素）。
\textbf{重要度 \xstarfull{4}}\quad\yiju{E4 2026 回忆版简答考"水击现象"；E1 slide33「第五章＝考试计算题重点」；E2「第五章同为主要计算考点」}\quad\nandu{中}

\begin{banfa}
\textbf{防护措施记成"三长两吸一泄压"}：延长关闭时间、延长（缩短）管路视情况、加大管径＝降流速；
空气包、调压塔＝吸收；安全阀／水锤消除器＝泄压。\\
\textbf{波速公式只记结构：}$c=\dfrac{\sqrt{K/\rho}}{\sqrt{1+DK/(\delta E)}}$。
考试若问"管壁弹性对水击的影响"，答"分母大于 1，$c$ 减小，$\Delta p$ 减小"。\\
\textbf{答题注意：}水击简答要写满\emph{定义＋原因＋危害＋措施}四块，2026 年原题就是简答，分值可观。
\end{banfa}
\end{kaodian}

\section{本章小结}

\begin{center}
\begin{tabularx}{\textwidth}{@{}lccX@{}}
\toprule
考点 & 重要度 & 难度 & 一句话抓住 \\
\midrule
水头损失两种形式与叠加 & \xstarfull{5} & 易 & 沿程＝黏性摩擦，局部＝漩涡耗能；$h_w=\sum h_f+\sum h_j$ \\
湿周、水力半径与当量直径 & \xstarfull{4} & 中 & $R_h=A/\chi$，$d_e=4R_h$；只能用于 $\Rey$ 与阻力折算 \\
雷诺实验与两种流态 & \xstarfull{5} & 易 & 直线／抖动／满管；层流 $h_f\propto v$，湍流 $h_f\propto v^{1.75\sim2}$ \\
雷诺数与流态判别 & \xstarfull{5} & 中 & $\Rey=vd/\nu$，下临界 2300、工程取 2000；$\Rey$＝惯性力／黏性力 \\
层流切应力与速度分布 & \xstarfull{4} & 中 & 切应力线性、速度抛物线；$v=u_{\max}/2$ \\
层流流量与 $\lambda=64/\Rey$ & \xstarfull{5} & 中 & $Q=\pi d^4\Delta p/(128\mu l)$；$\lambda$ 只与 $\Rey$ 有关 \\
湍流脉动与时均化 & \xstarfull{3} & 易 & 瞬时值＝时均值＋脉动值，脉动时均为零 \\
层流底层、水力光滑与粗糙 & \xstarfull{4} & 中 & $\delta>\Delta$ 光滑，$\delta<\Delta$ 粗糙；同一管随流量变 \\
尼古拉兹分区与 $\lambda$ 计算 & \xstarfull{5} & 难 & 层流 $64/\Rey$；光滑区 $0.3164/\Rey^{0.25}$；粗糙区只看 $\Delta/d$ \\
局部水头损失与突然扩大 & \xstarfull{4} & 中 & $h_j=\zeta v^2/2g$；突扩 $h_j=(v_1-v_2)^2/2g$ \\
典型 $\zeta$ 值与当量长度法 & \xstarfull{3} & 易 & 进口 0.5、出口 1.0；$h_j=\lambda(l_e/d)v^2/2g$ \\
薄壁小孔口自由出流 & \xstarfull{5} & 中 & $Q=\mu A\sqrt{2gH_0}$；$\varepsilon\approx0.64$、$\varphi\approx0.97$、$\mu\approx0.62$ \\
大孔口与淹没出流 & \xstarfull{3} & 中 & 大孔口 $e\ge H/10$；淹没出流 $H$＝上下游液面差 \\
圆柱形外管嘴与正常条件 & \xstarfull{5} & 中 & $\mu=0.82$；$H_0\le9\ \mathrm{m}$ 且 $l=(3\sim4)d$，缺一不可 \\
长管与短管的判别 & \xstarfull{5} & 中 & 不按长度分；长管忽略 $h_j$ 与 $v^2/2g$，短管必须计入 \\
简单长管三个基本问题 & \xstarfull{5} & 难 & 第一类直接算；第二、三类必须迭代 \\
串联与并联管路 & \xstarfull{5} & 难 & 串联流量相等、损失相加；并联损失相等、流量相加 \\
分支管路与管网 & \xstarfull{4} & 难 & 节点流量平衡，节点测压管水头唯一 \\
短管、虹吸管与水泵管路 & \xstarfull{4} & 难 & $Q=\mu A\sqrt{2gH_0}$；虹吸管最高处压强小于大气压 \\
水击现象、相长与分类 & \xstarfull{4} & 中 & $T=2L/c$；$t\le T$ 直接水击，$t>T$ 间接水击 \\
水击波速与防护 & \xstarfull{4} & 中 & $c=\sqrt{K/\rho}/\sqrt{1+DK/(\delta E)}$；延长关闭时间等 \\
\bottomrule
\end{tabularx}
\end{center}
